\documentclass[11pt,a4paper]{article}
% =====================================================================
%  Gemeinsame Vorlage der Arbeitsblätter "Binärbäume und Tiefensuche"
%  – Informatik 12, Kapitel 2. Kopie von ../../rekursive-datenstrukturen/
%  arbeitsblaetter/ab-vorlage.tex (gleiches Layout), Antwortplatz nach CLAUDE.md
%  immer eingeschaltet, dazu Stile für Binärbäume (baum, bk, ab, …).
%  Die Tiefensuche (../../tiefensuche/arbeitsblaetter/) bindet diese Datei ein.
%  Einbinden in jedem Blatt direkt nach \documentclass: % =====================================================================
%  Gemeinsame Vorlage der Arbeitsblätter "Binärbäume und Tiefensuche"
%  – Informatik 12, Kapitel 2. Kopie von ../../rekursive-datenstrukturen/
%  arbeitsblaetter/ab-vorlage.tex (gleiches Layout), Antwortplatz nach CLAUDE.md
%  immer eingeschaltet, dazu Stile für Binärbäume (baum, bk, ab, …).
%  Die Tiefensuche (../../tiefensuche/arbeitsblaetter/) bindet diese Datei ein.
%  Einbinden in jedem Blatt direkt nach \documentclass: % =====================================================================
%  Gemeinsame Vorlage der Arbeitsblätter "Binärbäume und Tiefensuche"
%  – Informatik 12, Kapitel 2. Kopie von ../../rekursive-datenstrukturen/
%  arbeitsblaetter/ab-vorlage.tex (gleiches Layout), Antwortplatz nach CLAUDE.md
%  immer eingeschaltet, dazu Stile für Binärbäume (baum, bk, ab, …).
%  Die Tiefensuche (../../tiefensuche/arbeitsblaetter/) bindet diese Datei ein.
%  Einbinden in jedem Blatt direkt nach \documentclass: \input{ab-vorlage}
%  Übersetzen mit XeLaTeX, LuaLaTeX, pdfLaTeX oder Tectonic.
%
%  Blatt:   xelatex ab-2-1-baeume.tex   (oder: tectonic ab-2-1-baeume.tex)
%  Lösung:  xelatex ab-2-1-baeume-lsg.tex
%           (Datei enthält nur: \def\mitloesung{}\input{ab-2-1-baeume};
%            füllt alle Lücken rot aus; siehe \lsg, \lsglinien, \lsgoder,
%            \lsgtext, lsgbild, \kreuz)
% =====================================================================
\usepackage[a4paper,left=16mm,right=16mm,top=23mm,bottom=13mm,headheight=12mm,headsep=3mm,footskip=6mm]{geometry}
\usepackage{iftex}
\ifPDFTeX
  \usepackage[T1]{fontenc}
  \usepackage[utf8]{inputenc}
  \usepackage{helvet}
  \usepackage[scaled=0.86]{beramono}
  \renewcommand{\familydefault}{\sfdefault}
\else
  \usepackage{fontspec}
  \setmainfont{texgyreheros}[Extension=.otf,UprightFont=*-regular,BoldFont=*-bold,ItalicFont=*-italic,BoldItalicFont=*-bolditalic]
  \setmonofont{DejaVuSansMono}[Extension=.ttf,UprightFont=*,BoldFont=*-Bold,ItalicFont=*-Oblique,BoldItalicFont=*-BoldOblique,Scale=0.86]
\fi
\usepackage[ngerman,shorthands=off]{babel}
\usepackage{amssymb}
\usepackage{xcolor,graphicx,tabularx,array,enumitem,fancyhdr,tikz,listings,multicol,calc,etoolbox,pifont}
\usepackage[most]{tcolorbox}
\usetikzlibrary{arrows.meta,positioning,calc,decorations.pathreplacing,shapes.multipart}
\usepackage[hidelinks]{hyperref}
\setlength{\parindent}{0pt}
\setlength{\parskip}{3pt}
\setlist{nosep,leftmargin=*}

% ---------- Lösungsschalter ----------
\newif\ifloesung
\ifdefined\mitloesung\loesungtrue\fi
\newcommand{\lsgfarbe}{\color{red!75!black}}

% ---------- Kopf- und Fußzeile (einheitliche Form) ----------
\newcommand{\lehrkraft}{Hau}
\newcommand{\themenbereich}{2. Binärbäume und Tiefensuche}
\newcommand{\abnummer}{}
\pagestyle{fancy}\fancyhf{}
\renewcommand{\headrulewidth}{0.4pt}
\fancyhead[L]{\small Klasse 12\\Datum:}
\fancyhead[C]{\small Themenbereich:\\\themenbereich}
\fancyhead[R]{\small \lehrkraft\ifloesung\\{\lsgfarbe\bfseries Lösung}\fi}
\fancyfoot[C]{\scriptsize Informatik 12\,·\,Blatt \abnummer\ifloesung\ (Lösung)\fi\,·\,Seite \thepage}
\newcommand{\abtitel}[2]{\renewcommand{\abnummer}{#1}\begin{center}\Large\bfseries #1\quad\underline{#2}\end{center}\vspace{-2mm}}

% ---------- Boxen ----------
\newenvironment{aufgabe}[2][]{\begin{tcolorbox}[enhanced,breakable,colback=white,colframe=black,boxrule=0.7pt,arc=3mm,left=2mm,right=2mm,top=1mm,bottom=1.2mm,before skip=2mm,after skip=2mm]\textbf{Aufgabe #2}\ifstrempty{#1}{}{ \textit{#1}}\par\vspace{0.3mm}}{\end{tcolorbox}}
\newtcolorbox{merke}[1][Merke]{enhanced,breakable,colback=green!5,colframe=green!40!black,boxrule=0.8pt,arc=2mm,left=2mm,right=2mm,top=1mm,bottom=1.2mm,before skip=2mm,after skip=2mm,title={#1},fonttitle=\bfseries,coltitle=white,toptitle=0.3mm,bottomtitle=0.3mm}
\newtcolorbox{hinweis}{enhanced,breakable,colback=yellow!12,colframe=orange!70!black,boxrule=0.6pt,arc=2mm,left=2mm,right=2mm,top=1mm,bottom=1mm,before skip=2mm,after skip=2mm}
\newcommand{\web}[1]{{\small\color{gray}\textit{Website: #1}}}

% ---------- Lücken, Linien, Karos ----------
\newcommand{\luecke}[1][3.5cm]{\rule[-0.5ex]{#1}{0.4pt}}
\newcommand{\linien}[1]{\par\vspace{1.5mm}\foreach \i in {1,...,#1}{\noindent\rule[0pt]{\linewidth}{0.3pt}\par\vspace{5.5mm}}\vspace{-3mm}}
\newcommand{\kariert}[1]{\par\vspace{1.5mm}\noindent\begin{tikzpicture}\draw[step=5mm,gray!55,thin] (0,0) grid (\linewidth-1pt,#1*5mm);\end{tikzpicture}\par}
\newcommand{\karobox}[2]{\begin{tikzpicture}\draw[step=5mm,gray!55,thin] (0,0) grid (#1,#2);\end{tikzpicture}}
\newcommand{\code}[1]{\texttt{#1}}
% Schritt-Nummern wie im Code und auf der Website: \nr{1} = ①
\newcommand{\nr}[1]{\ding{\numexpr171+#1\relax}}

% Lücke mit Lösung: \lsg[Mindestbreite]{Lösungstext}
%  Die Lösung steht (ohne eigene Höhe/Tiefe) auf derselben Linie wie im Blatt: Das Layout verschiebt sich nicht,
%  solange der Text nicht breiter als die Mindestbreite ist.
\newlength{\lsgbreite}
\newcommand{\lsg}[2][3.5cm]{\leavevmode\ifloesung\settowidth{\lsgbreite}{\,#2\,}\ifdim\lsgbreite<#1\setlength{\lsgbreite}{#1}\fi\rlap{\smash{\raisebox{0.8pt}{\makebox[\lsgbreite][l]{\lsgfarbe\,#2}}}}\luecke[\lsgbreite]\else\luecke[#1]\fi}
% Schreiblinien mit Lösung: \lsglinien{Anzahl}{Lösungstext} (gleiche Höhe wie \linien:
%  je Linie \baselineskip + 5,5 mm + \parskip, ausgemessen)
\newlength{\lsglinienhoehe}
\newcommand{\lsglinien}[2]{\ifloesung\par\vspace{1.5mm}\setlength{\lsglinienhoehe}{\dimexpr #1\dimexpr\baselineskip+5.5mm+\parskip\relax-16.25pt\relax}\noindent\begin{minipage}[t][\lsglinienhoehe][t]{\linewidth}\lsgfarbe #2\end{minipage}\par\vspace{-1.5mm}\else\linien{#1}\fi}
% Beliebiger Inhalt: \lsgoder{nur in der Lösung}{nur im Blatt}
\newcommand{\lsgoder}[2]{\ifloesung\leavevmode#1\else#2\fi}
% Text nur in der Lösung sichtbar, belegt im Blatt aber denselben Platz (ohne Text im PDF)
\newcommand{\lsgtext}[1]{\leavevmode\ifloesung{\lsgfarbe#1}\else\phantom{#1}\fi}
% Ankreuzfeld: \kreuz{1} = angekreuzt (nur in der Lösung), \kreuz{0} = leer
\newcommand{\kreuz}[1]{\leavevmode\ifloesung\ifnum#1=1 {\lsgfarbe$\boxtimes$}\else$\square$\fi\else$\square$\fi}
% Tabellenzeile mit fester Höhe für Lösungseinträge (2 Zeilen \small)
\newcommand{\zeile}{\rule[-6mm]{0pt}{9.5mm}}

% ---------- Verkettete Strukturen zeichnen ----------
% Objekt mit Kopf und Attributen:  \node[obj] (k1) {k1: Knoten \nodepart{second} daten = Hans\\ nachfolger};
\tikzset{
  obj/.style={draw,rectangle split,rectangle split parts=2,rectangle split part fill={orange!25,orange!8},
              rounded corners=2pt,font=\scriptsize,inner sep=2pt,align=left},
  qobj/.style={obj,rectangle split part fill={teal!25,teal!8}},
  aobj/.style={obj,rectangle split part fill={gray!30,gray!8}},
  pobj/.style={draw,rounded corners=2pt,fill=gray!8,font=\scriptsize,inner sep=2pt,align=left},
  klasse/.style={draw,rectangle split,rectangle split parts=3,font=\scriptsize,align=left,fill=white},
  zeiger/.style={thick,-{Stealth[length=2.2mm]}},
  erbt/.style={thick,-{Triangle[open,length=2.6mm,width=2.4mm]}},
  implementiert/.style={thick,densely dashed,-{Triangle[open,length=2.6mm,width=2.4mm]}},
  lsgzeiger/.style={zeiger,red!75!black},
  nullref/.style={font=\scriptsize\ttfamily,inner sep=1pt},
  schritt/.style={circle,draw,fill=white,inner sep=0.6pt,font=\scriptsize\bfseries},
  lsgschritt/.style={schritt,draw=red!75!black,text=red!75!black},
  % Sequenzdiagramm (UML 2): Kopf eckig, synchron = gefüllte Spitze, asynchron = offene Spitze,
  %  Rückgabe und Erzeugung gestrichelt mit offener Spitze (kein hohles Dreieck: das wäre Vererbung)
  lebenslinie/.style={draw,sharp corners,fill=white,font=\scriptsize,inner sep=2pt},
  leben/.style={densely dashed,gray},
  aufruf/.style={thick,-{Stealth[length=2mm]}},
  asynchron/.style={thick,-{Straight Barb[length=2mm,width=2.4mm]}},
  rueckgabe/.style={thick,densely dashed,-{Straight Barb[length=2mm,width=2.4mm]}},
  erzeugung/.style={rueckgabe},
}
% Akteur (Strichmännchen) mit Namen darunter: \akteur{x}{y oben}{Name}; ca. 0,9 cm hoch
\newcommand{\akteur}[3]{\begin{scope}[shift={(#1,#2)},line width=0.6pt]
  \draw (0,-0.09cm) circle (0.09cm);
  \draw (0,-0.18cm) -- (0,-0.42cm) (-0.16cm,-0.26cm) -- (0.16cm,-0.26cm)
        (-0.13cm,-0.62cm) -- (0,-0.42cm) -- (0.13cm,-0.62cm);
  \node[font=\scriptsize,inner sep=1pt,anchor=north] at (0,-0.63cm) {#3};
\end{scope}}
% Aktivierungsbalken auf der Lebenslinie bei x von y1 bis y2: \aktiv{x}{y1}{y2}
\newcommand{\aktiv}[3]{\draw[fill=gray!15] ($(#1,#2)+(-0.09cm,0)$) rectangle ($(#1,#3)+(0.09cm,0)$);}
% Lösungszeichnung: im Blatt unsichtbar, belegt aber denselben Platz
\ifloesung
  \tikzset{lsgbild/.style={red!75!black}}
\else
  \tikzset{lsgbild/.style={draw opacity=0,fill opacity=0,text opacity=0}}
\fi

% ---------- Java-Listings ----------
\lstset{language=Java,basicstyle=\ttfamily\small,keywordstyle=\bfseries,commentstyle=\color{gray!80!black},stringstyle=\color{black},showstringspaces=false,columns=flexible,keepspaces=true,tabsize=3,frame=single,framerule=0.4pt,rulecolor=\color{black},xleftmargin=2mm,framexleftmargin=1mm,aboveskip=2mm,belowskip=2mm,breaklines=true,escapechar=§,
 morekeywords={interface,implements,abstract,extends},
 literate={ä}{{\"a}}1 {ö}{{\"o}}1 {ü}{{\"u}}1 {Ä}{{\"A}}1 {Ö}{{\"O}}1 {Ü}{{\"U}}1 {ß}{{\ss}}1 {–}{{--}}1 {„}{{\glqq}}1 {“}{{\grqq}}1 {→}{{$\to$}}1}
\lstnewenvironment{java}[1][]{\lstset{#1}}{}

% ---------- Antwortplatz nach CLAUDE.md (in Kapitel 2 immer an) ----------
\def\antwortplatzacht{}
% Makros wie informatik/10/java/arbeitsblaetter/ab-vorlage.tex:
%  \linien{n}/\lsglinien{n}{Text}  Schreiblinien 8 mm      \linienfeld{n}/\lsgfeld{n}{Code}
%  \luecke[l]/\lsg[l]{Text}         Lücke (Standard 6,75 cm, Fachbegriff 5,25 cm, Zahl 2,25 cm)
%  \restzeile/\lsgrest{Text}        Lücke bis Zeilenende    \lueckentext  Absatz mit Lücken
%  \schreibzeile                    Tabellenzelle 8 mm      \lsgtext{…}  nur in der Lösung
%  \kariert{n}/\lsgkariert{n}{TikZ}, \karobox{b}{h}/\lsgkarobox{b}{h}{TikZ}  5-mm-Karo, ganze Kästchen
%  javaluecke                       Listing mit 1,5-fachem Zeilenabstand für Lücken (§\luecke§)
\ifdefined\antwortplatzacht
\usepackage{xspace}
\newlength{\zeilenabstand}\setlength{\zeilenabstand}{8mm}
\newlength{\karo}\setlength{\karo}{5mm}
\tikzset{karo/.style={draw=gray!50,line width=0.3pt},schreiblinie/.style={draw=black!65,line width=0.35pt}}
\renewcommand{\lsg}[2][6.75cm]{\leavevmode\ifloesung\settowidth{\lsgbreite}{\,#2\,}%
  \ifdim\lsgbreite>#1\relax\rlap{\raisebox{0.8pt}{\resizebox{#1}{!}{\lsgfarbe\,#2\,}}}%
  \else\rlap{\raisebox{0.8pt}{\lsgfarbe\,#2}}\fi\fi\rule[-0.6ex]{#1}{0.4pt}\xspace}
\renewcommand{\luecke}[1][6.75cm]{\lsg[#1]{}}
\newcommand{\lsgrest}[1]{\leavevmode\unskip\hspace{1.5mm}\ifloesung\rlap{\raisebox{0.8pt}{\lsgfarbe\,#1}}\fi\leaders\hbox{\rule[-0.6ex]{0.5pt}{0.4pt}}\hfill\null\par}
\newcommand{\restzeile}{\lsgrest{}}
\renewcommand{\lsgtext}[1]{\ifloesung\leavevmode{\lsgfarbe#1}\fi}
\newcommand{\lueckentext}{\linespread{1.5}\selectfont\raggedright}
\newcommand{\schreibzeile}{\rule[-2.8mm]{0pt}{8mm}}
\renewcommand{\lsglinien}[2]{\par\vspace{0.5mm}\noindent\begin{tikzpicture}
  \path (0,0) (\linewidth-0.4pt,-#1\zeilenabstand-1.5mm);
  \foreach \i in {1,...,#1}{\draw[schreiblinie] (0.2pt,-\i\zeilenabstand) -- (\linewidth-0.2pt,-\i\zeilenabstand);}
  \ifloesung\node[overlay,anchor=north west,inner sep=0pt] at (1mm,-\zeilenabstand+1pt+5mm) {\parbox[t]{\linewidth-2mm}{\lsgfarbe\raggedright\setlength{\baselineskip}{\zeilenabstand}\rule{0pt}{5mm}#2}};\fi
  \end{tikzpicture}\par}
\renewcommand{\linien}[1]{\lsglinien{#1}{}}
\newcommand{\lsgfeld}[2]{\noindent\begin{tikzpicture}
  \draw[line width=0.4pt] (0.2pt,0) rectangle (\linewidth-0.2pt,-#1\zeilenabstand-3mm);
  \foreach \i in {1,...,#1}{\draw[schreiblinie] (2mm,-\i\zeilenabstand) -- (\linewidth-2mm,-\i\zeilenabstand);}
  \ifloesung\node[overlay,anchor=north west,inner sep=0pt] at (2.5mm,-\zeilenabstand+1pt+5mm) {\parbox[t]{\linewidth-5mm}{\lsgfarbe\ttfamily\small\setlength{\baselineskip}{\zeilenabstand}\rule{0pt}{5mm}#2}};\fi
  \end{tikzpicture}}
\newcommand{\linienfeld}[1]{\lsgfeld{#1}{}}
\newcommand{\karoraster}[3]{\begin{tikzpicture}\draw[karo] (0,0) grid[step=\karo] (#1\karo,#2\karo);
  \ifloesung\begin{scope}[overlay,shift={(0,#2\karo)},red!75!black,text=red!75!black,line width=0.8pt]#3\end{scope}\fi\end{tikzpicture}}
\newcommand{\lsgkariert}[2]{\par\vspace{1.5mm}\noindent\pgfmathtruncatemacro{\karoX}{floor(\linewidth/\karo)}%
  \makebox[\linewidth]{\karoraster{\karoX}{#1}{#2}}\par\vspace{1mm}}
\renewcommand{\kariert}[1]{\lsgkariert{#1}{}}
\newcommand{\lsgkarobox}[3]{\pgfmathtruncatemacro{\karoX}{round((#1)/\karo)}\pgfmathtruncatemacro{\karoY}{round((#2)/\karo)}\karoraster{\karoX}{\karoY}{#3}}
\renewcommand{\karobox}[2]{\lsgkarobox{#1}{#2}{}}
\lstnewenvironment{javaluecke}[1][]{\lstset{basicstyle=\linespread{1.5}\small\ttfamily,#1}}{}
\fi

% ---------- Binärbäume ----------
% Knoten (bk), Abschluss (ab), markierter Knoten (bkm), Lösungsknoten (lsgbk, im Blatt unsichtbar)
\tikzset{
  baum/.style={level distance=9mm,level 1/.style={sibling distance=36mm},level 2/.style={sibling distance=18mm},
               level 3/.style={sibling distance=9mm},level 4/.style={sibling distance=6mm},
               edge from parent/.style={draw,thick,gray!60!black}},
  bk/.style={circle,draw=orange!70!black,fill=orange!22,minimum size=7.6mm,inner sep=0pt,font=\scriptsize,thick},
  bkm/.style={bk,fill=yellow!40,draw=orange!90!black,line width=1.2pt},
  ab/.style={rectangle,draw=gray!70,fill=gray!25,minimum size=2.2mm,inner sep=0pt},
  hoehe/.style={font=\scriptsize\bfseries,inner sep=1pt},
}
% Das Spielerverzeichnis (Mika, Emil, Tara, Bo, Ines, Pia, Zoe, Dana, Kai, Ruby, Vik); Argument: TikZ-Optionen
\newcommand{\verzeichnisbaum}[1][]{\begin{tikzpicture}[baum,#1]
\node[bk] (Mika) {Mika}
  child {node[bk] (Emil) {Emil}
    child {node[bk] (Bo) {Bo} child[missing] child {node[bk] (Dana) {Dana}}}
    child {node[bk] (Ines) {Ines} child[missing] child {node[bk] (Kai) {Kai}}}}
  child {node[bk] (Tara) {Tara}
    child {node[bk] (Pia) {Pia} child[missing] child {node[bk] (Ruby) {Ruby}}}
    child {node[bk] (Zoe) {Zoe} child {node[bk] (Vik) {Vik}} child[missing]}};
\end{tikzpicture}}

%  Übersetzen mit XeLaTeX, LuaLaTeX, pdfLaTeX oder Tectonic.
%
%  Blatt:   xelatex ab-2-1-baeume.tex   (oder: tectonic ab-2-1-baeume.tex)
%  Lösung:  xelatex ab-2-1-baeume-lsg.tex
%           (Datei enthält nur: \def\mitloesung{}\documentclass[11pt,a4paper]{article}
\input{ab-vorlage}
\begin{document}
\abtitel{2.1}{Bäume und geordnete Binärbäume}

Das Online-Spiel \textbf{Sternenstaub} hat zwei Millionen Accounts. Beim Login sucht der Server den Namen bisher in einer Liste. Du baust ein schnelleres Spielerverzeichnis. \web{Station 2.1}

\begin{aufgabe}[Namen raten]{1}
\begin{lueckentext}
\textbf{a)} Ich habe \lsg[2.25cm]{z.\,B. 8} Fragen gebraucht. Meine Strategie:\lsgrest{immer nach dem Namen in der Mitte fragen.}
\textbf{b)} Jede Frage nach der Mitte \lsg[5.25cm]{halbiert} den Bereich. Für 1000 Namen reichen deshalb höchstens \lsg[2.25cm]{10} Fragen, denn $2^{10} =$ \lsg[2.25cm]{1024}. Bei 2 Millionen Accounts sind es höchstens \lsg[2.25cm]{21} Fragen, eine Liste braucht bis zu \lsg[3.6cm]{2\,000\,000} Vergleiche.
\end{lueckentext}
\end{aufgabe}

\begin{merke}[Merke: Bäume]
\begin{lueckentext}
Ein Baum besteht aus \lsg[5.25cm]{Knoten}, die durch \lsg[5.25cm]{Kanten} verbunden sind. Der oberste Knoten heißt \lsg[5.25cm]{Wurzel}. Jeder andere Knoten hat genau einen \lsg[5.25cm]{Elternknoten}. Knoten ohne Kinder heißen \lsg[5.25cm]{Blätter}, alle anderen \lsg[5.25cm]{innere Knoten}. In einem \textbf{Binärbaum} hat jeder Knoten höchstens \lsg[2.25cm]{zwei} Kinder. Jeder Knoten bildet mit allem, was unter ihm hängt, einen \lsg[5.25cm]{Teilbaum}. Die Wurzel liegt auf Ebene 1. Die Anzahl der Ebenen ist die \lsg[5.25cm]{Höhe} des Baums.\par
\textbf{Rekursiv:} Ein Binärbaum ist entweder \lsg[5.25cm]{leer} oder er besteht aus einer Wurzel mit einem linken und einem rechten Teilbaum, die selbst wieder \lsg[5.25cm]{Binärbäume} sind.
\end{lueckentext}
\end{merke}

\begin{aufgabe}[Baum-Safari: das Spielerverzeichnis]{2}
\begin{minipage}[t]{0.5\linewidth}
\vspace{0pt}\centering
\verzeichnisbaum[scale=0.98,every node/.append style={transform shape}]
\end{minipage}\hfill
\begin{minipage}[t]{0.47\linewidth}
\vspace{0pt}\lueckentext
Wurzel: \lsg[3.6cm]{Mika}\par
Blätter: \lsgrest{Dana, Kai, Ruby, Vik}
Elternknoten von Kai: \lsg[2.6cm]{Ines}\par
Kinder von Tara: \lsg[4.2cm]{Pia, Zoe}\par
Rechter Teilbaum von Emil: \lsgrest{Ines, Kai}
Knoten auf Ebene 3: \lsgrest{Bo, Ines, Pia, Zoe}
Höhe des Baums: \lsg[2.25cm]{4}
\end{minipage}
\end{aufgabe}

\begin{merke}[Merke: Geordneter Binärbaum (Suchbaum)]
\begin{lueckentext}
Für \textbf{jeden} Knoten gilt: Alle Schlüssel im linken Teilbaum sind \lsg[5.25cm]{kleiner}, alle im rechten Teilbaum sind \lsg[5.25cm]{größer} als sein eigener Schlüssel. Jeder Schlüssel kommt nur einmal vor. (Hier: Schlüssel = Name, „kleiner“ = kommt im Alphabet früher.)\par
\end{lueckentext}
\begin{minipage}[t]{0.49\linewidth}
\textbf{Vorgehensweise Suchen} (ab der Wurzel):\\
\nr{1} gleich: gefunden.\\
\nr{2} kleiner: im linken Teilbaum weitersuchen.\\
\nr{3} größer: im rechten Teilbaum weitersuchen.\\
\nr{4} Teilbaum leer: Den Namen gibt es nicht.
\end{minipage}\hfill
\begin{minipage}[t]{0.49\linewidth}
\textbf{Vorgehensweise Einfügen} (ab der Wurzel):\\
\nr{1} gleich: Name vergeben, nicht einfügen.\\
\nr{2} kleiner: im linken Teilbaum einfügen.\\
\nr{3} größer: im rechten Teilbaum einfügen.\\
\nr{4} Teilbaum leer: Hier entsteht der neue Knoten.
\end{minipage}\par\vspace{1mm}
\begin{lueckentext}
Der neue Knoten ist immer ein \lsg[5.25cm]{Blatt}. Die Suche geht bei jedem Vergleich in nur \lsg[2.25cm]{einen} Teilbaum, sie braucht also höchstens so viele Vergleiche, wie der Baum \lsg[5.25cm]{Ebenen} hat.
\end{lueckentext}
\end{merke}

\newpage
\begin{aufgabe}[Such-Sprint: Suchwege im Verzeichnis (Baum aus Aufgabe 2)]{3}
{\small\renewcommand{\arraystretch}{1.1}
\begin{tabularx}{\linewidth}{|l|X|c|c|}\hline
gesucht & Suchweg (besuchte Knoten) & Vergleiche & gefunden? \\\hline
Ruby & \schreibzeile\lsgtext{Mika → Tara → Pia → Ruby} & \lsgtext{4} & \lsgtext{ja} \\\hline
Dana & \schreibzeile\lsgtext{Mika → Emil → Bo → Dana} & \lsgtext{4} & \lsgtext{ja} \\\hline
Olaf & \schreibzeile\lsgtext{Mika → Tara → Pia → (links unter Pia leer)} & \lsgtext{3} & \lsgtext{nein} \\\hline
\end{tabularx}}
\end{aufgabe}

\begin{aufgabe}[Baumschule]{4}
\textbf{a)} Zeichne den Baum zur Reihenfolge \code{Ole Kim Tom Ben Lia Sam Zara}.\hfill Höhe: \lsg[2.25cm]{3}\par\vspace{1mm}
\lsgkariert{6}{
  \begin{scope}[baum,level distance=9mm,level 1/.style={sibling distance=44mm},level 2/.style={sibling distance=20mm},every node/.append style={draw=red!75!black,fill=white}]
  \node[bk] at (8.5,-0.6) {Ole}
    child {node[bk] {Kim} child {node[bk] {Ben}} child {node[bk] {Lia}}}
    child {node[bk] {Tom} child {node[bk] {Sam}} child {node[bk] {Zara}}};
  \end{scope}}\par\vspace{1mm}
\textbf{b)} Skizziere den Baum für dieselben Namen alphabetisch.\hfill Höhe: \lsg[2.25cm]{7}\par\vspace{1mm}
\begin{minipage}[t]{0.52\linewidth}
\vspace{0pt}
\lsgkarobox{8.5cm}{4cm}{
  \foreach \n/\i in {Ben/0,Kim/1,Lia/2,Ole/3,Sam/4,Tom/5,Zara/6}{
    \node[bk,minimum size=6mm,font=\tiny,draw=red!75!black,fill=white] (z\i) at (0.6+1.2*\i,-0.45-0.5*\i) {\n};}
  \foreach \i [evaluate=\i as \j using int(\i+1)] in {0,...,5}{\draw[thick] (z\i) -- (z\j);}}
\end{minipage}\hfill
\begin{minipage}[t]{0.45\linewidth}
\vspace{0pt}\lueckentext
Werden die Namen sortiert eingefügt,\par\lsgrest{hängt jeder neue Name ganz rechts:}
\lsgrest{eine Kette, so langsam wie eine Liste}
\lsgrest{(entarteter Baum).}
\end{minipage}
\end{aufgabe}

\begin{aufgabe}[Ist das ein geordneter Binärbaum?]{5}
\begin{minipage}[t]{0.38\linewidth}
\vspace{0pt}\centering
\begin{tikzpicture}[baum,level distance=8mm,level 1/.style={sibling distance=30mm},level 2/.style={sibling distance=15mm},level 3/.style={sibling distance=8mm},every node/.append style={scale=0.9}]
\node[bk] {Mika}
  child {node[bk] {Emil} child {node[bk] {Bo}} child {node[bk] {Ines}}}
  child {node[bk] {Tara} child {node[bk] {Pia} child {node[bk] {Kai}} child[missing]} child {node[bk] {Zoe}}};
\end{tikzpicture}
\end{minipage}\hfill
\begin{minipage}[t]{0.6\linewidth}
\vspace{0pt}
Jemand hat Kai von Hand links unter Pia eingetragen, „weil Kai vor Pia kommt“.\par\vspace{1mm}
\kreuz{0} ja, geordnet\qquad \kreuz{1} nein, nicht geordnet, weil:\par
\lsglinien{3}{Kai steht im rechten Teilbaum von Mika, kommt aber vor Mika. Die Suche nach Kai biegt bei Mika links ab und findet ihn nie.}
\vspace{2mm}\begin{lueckentext}
Kai gehört \lsgrest{rechts unter Ines.}
\end{lueckentext}
\end{minipage}
\end{aufgabe}
\end{document}
;
%            füllt alle Lücken rot aus; siehe \lsg, \lsglinien, \lsgoder,
%            \lsgtext, lsgbild, \kreuz)
% =====================================================================
\usepackage[a4paper,left=16mm,right=16mm,top=23mm,bottom=13mm,headheight=12mm,headsep=3mm,footskip=6mm]{geometry}
\usepackage{iftex}
\ifPDFTeX
  \usepackage[T1]{fontenc}
  \usepackage[utf8]{inputenc}
  \usepackage{helvet}
  \usepackage[scaled=0.86]{beramono}
  \renewcommand{\familydefault}{\sfdefault}
\else
  \usepackage{fontspec}
  \setmainfont{texgyreheros}[Extension=.otf,UprightFont=*-regular,BoldFont=*-bold,ItalicFont=*-italic,BoldItalicFont=*-bolditalic]
  \setmonofont{DejaVuSansMono}[Extension=.ttf,UprightFont=*,BoldFont=*-Bold,ItalicFont=*-Oblique,BoldItalicFont=*-BoldOblique,Scale=0.86]
\fi
\usepackage[ngerman,shorthands=off]{babel}
\usepackage{amssymb}
\usepackage{xcolor,graphicx,tabularx,array,enumitem,fancyhdr,tikz,listings,multicol,calc,etoolbox,pifont}
\usepackage[most]{tcolorbox}
\usetikzlibrary{arrows.meta,positioning,calc,decorations.pathreplacing,shapes.multipart}
\usepackage[hidelinks]{hyperref}
\setlength{\parindent}{0pt}
\setlength{\parskip}{3pt}
\setlist{nosep,leftmargin=*}

% ---------- Lösungsschalter ----------
\newif\ifloesung
\ifdefined\mitloesung\loesungtrue\fi
\newcommand{\lsgfarbe}{\color{red!75!black}}

% ---------- Kopf- und Fußzeile (einheitliche Form) ----------
\newcommand{\lehrkraft}{Hau}
\newcommand{\themenbereich}{2. Binärbäume und Tiefensuche}
\newcommand{\abnummer}{}
\pagestyle{fancy}\fancyhf{}
\renewcommand{\headrulewidth}{0.4pt}
\fancyhead[L]{\small Klasse 12\\Datum:}
\fancyhead[C]{\small Themenbereich:\\\themenbereich}
\fancyhead[R]{\small \lehrkraft\ifloesung\\{\lsgfarbe\bfseries Lösung}\fi}
\fancyfoot[C]{\scriptsize Informatik 12\,·\,Blatt \abnummer\ifloesung\ (Lösung)\fi\,·\,Seite \thepage}
\newcommand{\abtitel}[2]{\renewcommand{\abnummer}{#1}\begin{center}\Large\bfseries #1\quad\underline{#2}\end{center}\vspace{-2mm}}

% ---------- Boxen ----------
\newenvironment{aufgabe}[2][]{\begin{tcolorbox}[enhanced,breakable,colback=white,colframe=black,boxrule=0.7pt,arc=3mm,left=2mm,right=2mm,top=1mm,bottom=1.2mm,before skip=2mm,after skip=2mm]\textbf{Aufgabe #2}\ifstrempty{#1}{}{ \textit{#1}}\par\vspace{0.3mm}}{\end{tcolorbox}}
\newtcolorbox{merke}[1][Merke]{enhanced,breakable,colback=green!5,colframe=green!40!black,boxrule=0.8pt,arc=2mm,left=2mm,right=2mm,top=1mm,bottom=1.2mm,before skip=2mm,after skip=2mm,title={#1},fonttitle=\bfseries,coltitle=white,toptitle=0.3mm,bottomtitle=0.3mm}
\newtcolorbox{hinweis}{enhanced,breakable,colback=yellow!12,colframe=orange!70!black,boxrule=0.6pt,arc=2mm,left=2mm,right=2mm,top=1mm,bottom=1mm,before skip=2mm,after skip=2mm}
\newcommand{\web}[1]{{\small\color{gray}\textit{Website: #1}}}

% ---------- Lücken, Linien, Karos ----------
\newcommand{\luecke}[1][3.5cm]{\rule[-0.5ex]{#1}{0.4pt}}
\newcommand{\linien}[1]{\par\vspace{1.5mm}\foreach \i in {1,...,#1}{\noindent\rule[0pt]{\linewidth}{0.3pt}\par\vspace{5.5mm}}\vspace{-3mm}}
\newcommand{\kariert}[1]{\par\vspace{1.5mm}\noindent\begin{tikzpicture}\draw[step=5mm,gray!55,thin] (0,0) grid (\linewidth-1pt,#1*5mm);\end{tikzpicture}\par}
\newcommand{\karobox}[2]{\begin{tikzpicture}\draw[step=5mm,gray!55,thin] (0,0) grid (#1,#2);\end{tikzpicture}}
\newcommand{\code}[1]{\texttt{#1}}
% Schritt-Nummern wie im Code und auf der Website: \nr{1} = ①
\newcommand{\nr}[1]{\ding{\numexpr171+#1\relax}}

% Lücke mit Lösung: \lsg[Mindestbreite]{Lösungstext}
%  Die Lösung steht (ohne eigene Höhe/Tiefe) auf derselben Linie wie im Blatt: Das Layout verschiebt sich nicht,
%  solange der Text nicht breiter als die Mindestbreite ist.
\newlength{\lsgbreite}
\newcommand{\lsg}[2][3.5cm]{\leavevmode\ifloesung\settowidth{\lsgbreite}{\,#2\,}\ifdim\lsgbreite<#1\setlength{\lsgbreite}{#1}\fi\rlap{\smash{\raisebox{0.8pt}{\makebox[\lsgbreite][l]{\lsgfarbe\,#2}}}}\luecke[\lsgbreite]\else\luecke[#1]\fi}
% Schreiblinien mit Lösung: \lsglinien{Anzahl}{Lösungstext} (gleiche Höhe wie \linien:
%  je Linie \baselineskip + 5,5 mm + \parskip, ausgemessen)
\newlength{\lsglinienhoehe}
\newcommand{\lsglinien}[2]{\ifloesung\par\vspace{1.5mm}\setlength{\lsglinienhoehe}{\dimexpr #1\dimexpr\baselineskip+5.5mm+\parskip\relax-16.25pt\relax}\noindent\begin{minipage}[t][\lsglinienhoehe][t]{\linewidth}\lsgfarbe #2\end{minipage}\par\vspace{-1.5mm}\else\linien{#1}\fi}
% Beliebiger Inhalt: \lsgoder{nur in der Lösung}{nur im Blatt}
\newcommand{\lsgoder}[2]{\ifloesung\leavevmode#1\else#2\fi}
% Text nur in der Lösung sichtbar, belegt im Blatt aber denselben Platz (ohne Text im PDF)
\newcommand{\lsgtext}[1]{\leavevmode\ifloesung{\lsgfarbe#1}\else\phantom{#1}\fi}
% Ankreuzfeld: \kreuz{1} = angekreuzt (nur in der Lösung), \kreuz{0} = leer
\newcommand{\kreuz}[1]{\leavevmode\ifloesung\ifnum#1=1 {\lsgfarbe$\boxtimes$}\else$\square$\fi\else$\square$\fi}
% Tabellenzeile mit fester Höhe für Lösungseinträge (2 Zeilen \small)
\newcommand{\zeile}{\rule[-6mm]{0pt}{9.5mm}}

% ---------- Verkettete Strukturen zeichnen ----------
% Objekt mit Kopf und Attributen:  \node[obj] (k1) {k1: Knoten \nodepart{second} daten = Hans\\ nachfolger};
\tikzset{
  obj/.style={draw,rectangle split,rectangle split parts=2,rectangle split part fill={orange!25,orange!8},
              rounded corners=2pt,font=\scriptsize,inner sep=2pt,align=left},
  qobj/.style={obj,rectangle split part fill={teal!25,teal!8}},
  aobj/.style={obj,rectangle split part fill={gray!30,gray!8}},
  pobj/.style={draw,rounded corners=2pt,fill=gray!8,font=\scriptsize,inner sep=2pt,align=left},
  klasse/.style={draw,rectangle split,rectangle split parts=3,font=\scriptsize,align=left,fill=white},
  zeiger/.style={thick,-{Stealth[length=2.2mm]}},
  erbt/.style={thick,-{Triangle[open,length=2.6mm,width=2.4mm]}},
  implementiert/.style={thick,densely dashed,-{Triangle[open,length=2.6mm,width=2.4mm]}},
  lsgzeiger/.style={zeiger,red!75!black},
  nullref/.style={font=\scriptsize\ttfamily,inner sep=1pt},
  schritt/.style={circle,draw,fill=white,inner sep=0.6pt,font=\scriptsize\bfseries},
  lsgschritt/.style={schritt,draw=red!75!black,text=red!75!black},
  % Sequenzdiagramm (UML 2): Kopf eckig, synchron = gefüllte Spitze, asynchron = offene Spitze,
  %  Rückgabe und Erzeugung gestrichelt mit offener Spitze (kein hohles Dreieck: das wäre Vererbung)
  lebenslinie/.style={draw,sharp corners,fill=white,font=\scriptsize,inner sep=2pt},
  leben/.style={densely dashed,gray},
  aufruf/.style={thick,-{Stealth[length=2mm]}},
  asynchron/.style={thick,-{Straight Barb[length=2mm,width=2.4mm]}},
  rueckgabe/.style={thick,densely dashed,-{Straight Barb[length=2mm,width=2.4mm]}},
  erzeugung/.style={rueckgabe},
}
% Akteur (Strichmännchen) mit Namen darunter: \akteur{x}{y oben}{Name}; ca. 0,9 cm hoch
\newcommand{\akteur}[3]{\begin{scope}[shift={(#1,#2)},line width=0.6pt]
  \draw (0,-0.09cm) circle (0.09cm);
  \draw (0,-0.18cm) -- (0,-0.42cm) (-0.16cm,-0.26cm) -- (0.16cm,-0.26cm)
        (-0.13cm,-0.62cm) -- (0,-0.42cm) -- (0.13cm,-0.62cm);
  \node[font=\scriptsize,inner sep=1pt,anchor=north] at (0,-0.63cm) {#3};
\end{scope}}
% Aktivierungsbalken auf der Lebenslinie bei x von y1 bis y2: \aktiv{x}{y1}{y2}
\newcommand{\aktiv}[3]{\draw[fill=gray!15] ($(#1,#2)+(-0.09cm,0)$) rectangle ($(#1,#3)+(0.09cm,0)$);}
% Lösungszeichnung: im Blatt unsichtbar, belegt aber denselben Platz
\ifloesung
  \tikzset{lsgbild/.style={red!75!black}}
\else
  \tikzset{lsgbild/.style={draw opacity=0,fill opacity=0,text opacity=0}}
\fi

% ---------- Java-Listings ----------
\lstset{language=Java,basicstyle=\ttfamily\small,keywordstyle=\bfseries,commentstyle=\color{gray!80!black},stringstyle=\color{black},showstringspaces=false,columns=flexible,keepspaces=true,tabsize=3,frame=single,framerule=0.4pt,rulecolor=\color{black},xleftmargin=2mm,framexleftmargin=1mm,aboveskip=2mm,belowskip=2mm,breaklines=true,escapechar=§,
 morekeywords={interface,implements,abstract,extends},
 literate={ä}{{\"a}}1 {ö}{{\"o}}1 {ü}{{\"u}}1 {Ä}{{\"A}}1 {Ö}{{\"O}}1 {Ü}{{\"U}}1 {ß}{{\ss}}1 {–}{{--}}1 {„}{{\glqq}}1 {“}{{\grqq}}1 {→}{{$\to$}}1}
\lstnewenvironment{java}[1][]{\lstset{#1}}{}

% ---------- Antwortplatz nach CLAUDE.md (in Kapitel 2 immer an) ----------
\def\antwortplatzacht{}
% Makros wie informatik/10/java/arbeitsblaetter/ab-vorlage.tex:
%  \linien{n}/\lsglinien{n}{Text}  Schreiblinien 8 mm      \linienfeld{n}/\lsgfeld{n}{Code}
%  \luecke[l]/\lsg[l]{Text}         Lücke (Standard 6,75 cm, Fachbegriff 5,25 cm, Zahl 2,25 cm)
%  \restzeile/\lsgrest{Text}        Lücke bis Zeilenende    \lueckentext  Absatz mit Lücken
%  \schreibzeile                    Tabellenzelle 8 mm      \lsgtext{…}  nur in der Lösung
%  \kariert{n}/\lsgkariert{n}{TikZ}, \karobox{b}{h}/\lsgkarobox{b}{h}{TikZ}  5-mm-Karo, ganze Kästchen
%  javaluecke                       Listing mit 1,5-fachem Zeilenabstand für Lücken (§\luecke§)
\ifdefined\antwortplatzacht
\usepackage{xspace}
\newlength{\zeilenabstand}\setlength{\zeilenabstand}{8mm}
\newlength{\karo}\setlength{\karo}{5mm}
\tikzset{karo/.style={draw=gray!50,line width=0.3pt},schreiblinie/.style={draw=black!65,line width=0.35pt}}
\renewcommand{\lsg}[2][6.75cm]{\leavevmode\ifloesung\settowidth{\lsgbreite}{\,#2\,}%
  \ifdim\lsgbreite>#1\relax\rlap{\raisebox{0.8pt}{\resizebox{#1}{!}{\lsgfarbe\,#2\,}}}%
  \else\rlap{\raisebox{0.8pt}{\lsgfarbe\,#2}}\fi\fi\rule[-0.6ex]{#1}{0.4pt}\xspace}
\renewcommand{\luecke}[1][6.75cm]{\lsg[#1]{}}
\newcommand{\lsgrest}[1]{\leavevmode\unskip\hspace{1.5mm}\ifloesung\rlap{\raisebox{0.8pt}{\lsgfarbe\,#1}}\fi\leaders\hbox{\rule[-0.6ex]{0.5pt}{0.4pt}}\hfill\null\par}
\newcommand{\restzeile}{\lsgrest{}}
\renewcommand{\lsgtext}[1]{\ifloesung\leavevmode{\lsgfarbe#1}\fi}
\newcommand{\lueckentext}{\linespread{1.5}\selectfont\raggedright}
\newcommand{\schreibzeile}{\rule[-2.8mm]{0pt}{8mm}}
\renewcommand{\lsglinien}[2]{\par\vspace{0.5mm}\noindent\begin{tikzpicture}
  \path (0,0) (\linewidth-0.4pt,-#1\zeilenabstand-1.5mm);
  \foreach \i in {1,...,#1}{\draw[schreiblinie] (0.2pt,-\i\zeilenabstand) -- (\linewidth-0.2pt,-\i\zeilenabstand);}
  \ifloesung\node[overlay,anchor=north west,inner sep=0pt] at (1mm,-\zeilenabstand+1pt+5mm) {\parbox[t]{\linewidth-2mm}{\lsgfarbe\raggedright\setlength{\baselineskip}{\zeilenabstand}\rule{0pt}{5mm}#2}};\fi
  \end{tikzpicture}\par}
\renewcommand{\linien}[1]{\lsglinien{#1}{}}
\newcommand{\lsgfeld}[2]{\noindent\begin{tikzpicture}
  \draw[line width=0.4pt] (0.2pt,0) rectangle (\linewidth-0.2pt,-#1\zeilenabstand-3mm);
  \foreach \i in {1,...,#1}{\draw[schreiblinie] (2mm,-\i\zeilenabstand) -- (\linewidth-2mm,-\i\zeilenabstand);}
  \ifloesung\node[overlay,anchor=north west,inner sep=0pt] at (2.5mm,-\zeilenabstand+1pt+5mm) {\parbox[t]{\linewidth-5mm}{\lsgfarbe\ttfamily\small\setlength{\baselineskip}{\zeilenabstand}\rule{0pt}{5mm}#2}};\fi
  \end{tikzpicture}}
\newcommand{\linienfeld}[1]{\lsgfeld{#1}{}}
\newcommand{\karoraster}[3]{\begin{tikzpicture}\draw[karo] (0,0) grid[step=\karo] (#1\karo,#2\karo);
  \ifloesung\begin{scope}[overlay,shift={(0,#2\karo)},red!75!black,text=red!75!black,line width=0.8pt]#3\end{scope}\fi\end{tikzpicture}}
\newcommand{\lsgkariert}[2]{\par\vspace{1.5mm}\noindent\pgfmathtruncatemacro{\karoX}{floor(\linewidth/\karo)}%
  \makebox[\linewidth]{\karoraster{\karoX}{#1}{#2}}\par\vspace{1mm}}
\renewcommand{\kariert}[1]{\lsgkariert{#1}{}}
\newcommand{\lsgkarobox}[3]{\pgfmathtruncatemacro{\karoX}{round((#1)/\karo)}\pgfmathtruncatemacro{\karoY}{round((#2)/\karo)}\karoraster{\karoX}{\karoY}{#3}}
\renewcommand{\karobox}[2]{\lsgkarobox{#1}{#2}{}}
\lstnewenvironment{javaluecke}[1][]{\lstset{basicstyle=\linespread{1.5}\small\ttfamily,#1}}{}
\fi

% ---------- Binärbäume ----------
% Knoten (bk), Abschluss (ab), markierter Knoten (bkm), Lösungsknoten (lsgbk, im Blatt unsichtbar)
\tikzset{
  baum/.style={level distance=9mm,level 1/.style={sibling distance=36mm},level 2/.style={sibling distance=18mm},
               level 3/.style={sibling distance=9mm},level 4/.style={sibling distance=6mm},
               edge from parent/.style={draw,thick,gray!60!black}},
  bk/.style={circle,draw=orange!70!black,fill=orange!22,minimum size=7.6mm,inner sep=0pt,font=\scriptsize,thick},
  bkm/.style={bk,fill=yellow!40,draw=orange!90!black,line width=1.2pt},
  ab/.style={rectangle,draw=gray!70,fill=gray!25,minimum size=2.2mm,inner sep=0pt},
  hoehe/.style={font=\scriptsize\bfseries,inner sep=1pt},
}
% Das Spielerverzeichnis (Mika, Emil, Tara, Bo, Ines, Pia, Zoe, Dana, Kai, Ruby, Vik); Argument: TikZ-Optionen
\newcommand{\verzeichnisbaum}[1][]{\begin{tikzpicture}[baum,#1]
\node[bk] (Mika) {Mika}
  child {node[bk] (Emil) {Emil}
    child {node[bk] (Bo) {Bo} child[missing] child {node[bk] (Dana) {Dana}}}
    child {node[bk] (Ines) {Ines} child[missing] child {node[bk] (Kai) {Kai}}}}
  child {node[bk] (Tara) {Tara}
    child {node[bk] (Pia) {Pia} child[missing] child {node[bk] (Ruby) {Ruby}}}
    child {node[bk] (Zoe) {Zoe} child {node[bk] (Vik) {Vik}} child[missing]}};
\end{tikzpicture}}

%  Übersetzen mit XeLaTeX, LuaLaTeX, pdfLaTeX oder Tectonic.
%
%  Blatt:   xelatex ab-2-1-baeume.tex   (oder: tectonic ab-2-1-baeume.tex)
%  Lösung:  xelatex ab-2-1-baeume-lsg.tex
%           (Datei enthält nur: \def\mitloesung{}\documentclass[11pt,a4paper]{article}
% =====================================================================
%  Gemeinsame Vorlage der Arbeitsblätter "Binärbäume und Tiefensuche"
%  – Informatik 12, Kapitel 2. Kopie von ../../rekursive-datenstrukturen/
%  arbeitsblaetter/ab-vorlage.tex (gleiches Layout), Antwortplatz nach CLAUDE.md
%  immer eingeschaltet, dazu Stile für Binärbäume (baum, bk, ab, …).
%  Die Tiefensuche (../../tiefensuche/arbeitsblaetter/) bindet diese Datei ein.
%  Einbinden in jedem Blatt direkt nach \documentclass: \input{ab-vorlage}
%  Übersetzen mit XeLaTeX, LuaLaTeX, pdfLaTeX oder Tectonic.
%
%  Blatt:   xelatex ab-2-1-baeume.tex   (oder: tectonic ab-2-1-baeume.tex)
%  Lösung:  xelatex ab-2-1-baeume-lsg.tex
%           (Datei enthält nur: \def\mitloesung{}\input{ab-2-1-baeume};
%            füllt alle Lücken rot aus; siehe \lsg, \lsglinien, \lsgoder,
%            \lsgtext, lsgbild, \kreuz)
% =====================================================================
\usepackage[a4paper,left=16mm,right=16mm,top=23mm,bottom=13mm,headheight=12mm,headsep=3mm,footskip=6mm]{geometry}
\usepackage{iftex}
\ifPDFTeX
  \usepackage[T1]{fontenc}
  \usepackage[utf8]{inputenc}
  \usepackage{helvet}
  \usepackage[scaled=0.86]{beramono}
  \renewcommand{\familydefault}{\sfdefault}
\else
  \usepackage{fontspec}
  \setmainfont{texgyreheros}[Extension=.otf,UprightFont=*-regular,BoldFont=*-bold,ItalicFont=*-italic,BoldItalicFont=*-bolditalic]
  \setmonofont{DejaVuSansMono}[Extension=.ttf,UprightFont=*,BoldFont=*-Bold,ItalicFont=*-Oblique,BoldItalicFont=*-BoldOblique,Scale=0.86]
\fi
\usepackage[ngerman,shorthands=off]{babel}
\usepackage{amssymb}
\usepackage{xcolor,graphicx,tabularx,array,enumitem,fancyhdr,tikz,listings,multicol,calc,etoolbox,pifont}
\usepackage[most]{tcolorbox}
\usetikzlibrary{arrows.meta,positioning,calc,decorations.pathreplacing,shapes.multipart}
\usepackage[hidelinks]{hyperref}
\setlength{\parindent}{0pt}
\setlength{\parskip}{3pt}
\setlist{nosep,leftmargin=*}

% ---------- Lösungsschalter ----------
\newif\ifloesung
\ifdefined\mitloesung\loesungtrue\fi
\newcommand{\lsgfarbe}{\color{red!75!black}}

% ---------- Kopf- und Fußzeile (einheitliche Form) ----------
\newcommand{\lehrkraft}{Hau}
\newcommand{\themenbereich}{2. Binärbäume und Tiefensuche}
\newcommand{\abnummer}{}
\pagestyle{fancy}\fancyhf{}
\renewcommand{\headrulewidth}{0.4pt}
\fancyhead[L]{\small Klasse 12\\Datum:}
\fancyhead[C]{\small Themenbereich:\\\themenbereich}
\fancyhead[R]{\small \lehrkraft\ifloesung\\{\lsgfarbe\bfseries Lösung}\fi}
\fancyfoot[C]{\scriptsize Informatik 12\,·\,Blatt \abnummer\ifloesung\ (Lösung)\fi\,·\,Seite \thepage}
\newcommand{\abtitel}[2]{\renewcommand{\abnummer}{#1}\begin{center}\Large\bfseries #1\quad\underline{#2}\end{center}\vspace{-2mm}}

% ---------- Boxen ----------
\newenvironment{aufgabe}[2][]{\begin{tcolorbox}[enhanced,breakable,colback=white,colframe=black,boxrule=0.7pt,arc=3mm,left=2mm,right=2mm,top=1mm,bottom=1.2mm,before skip=2mm,after skip=2mm]\textbf{Aufgabe #2}\ifstrempty{#1}{}{ \textit{#1}}\par\vspace{0.3mm}}{\end{tcolorbox}}
\newtcolorbox{merke}[1][Merke]{enhanced,breakable,colback=green!5,colframe=green!40!black,boxrule=0.8pt,arc=2mm,left=2mm,right=2mm,top=1mm,bottom=1.2mm,before skip=2mm,after skip=2mm,title={#1},fonttitle=\bfseries,coltitle=white,toptitle=0.3mm,bottomtitle=0.3mm}
\newtcolorbox{hinweis}{enhanced,breakable,colback=yellow!12,colframe=orange!70!black,boxrule=0.6pt,arc=2mm,left=2mm,right=2mm,top=1mm,bottom=1mm,before skip=2mm,after skip=2mm}
\newcommand{\web}[1]{{\small\color{gray}\textit{Website: #1}}}

% ---------- Lücken, Linien, Karos ----------
\newcommand{\luecke}[1][3.5cm]{\rule[-0.5ex]{#1}{0.4pt}}
\newcommand{\linien}[1]{\par\vspace{1.5mm}\foreach \i in {1,...,#1}{\noindent\rule[0pt]{\linewidth}{0.3pt}\par\vspace{5.5mm}}\vspace{-3mm}}
\newcommand{\kariert}[1]{\par\vspace{1.5mm}\noindent\begin{tikzpicture}\draw[step=5mm,gray!55,thin] (0,0) grid (\linewidth-1pt,#1*5mm);\end{tikzpicture}\par}
\newcommand{\karobox}[2]{\begin{tikzpicture}\draw[step=5mm,gray!55,thin] (0,0) grid (#1,#2);\end{tikzpicture}}
\newcommand{\code}[1]{\texttt{#1}}
% Schritt-Nummern wie im Code und auf der Website: \nr{1} = ①
\newcommand{\nr}[1]{\ding{\numexpr171+#1\relax}}

% Lücke mit Lösung: \lsg[Mindestbreite]{Lösungstext}
%  Die Lösung steht (ohne eigene Höhe/Tiefe) auf derselben Linie wie im Blatt: Das Layout verschiebt sich nicht,
%  solange der Text nicht breiter als die Mindestbreite ist.
\newlength{\lsgbreite}
\newcommand{\lsg}[2][3.5cm]{\leavevmode\ifloesung\settowidth{\lsgbreite}{\,#2\,}\ifdim\lsgbreite<#1\setlength{\lsgbreite}{#1}\fi\rlap{\smash{\raisebox{0.8pt}{\makebox[\lsgbreite][l]{\lsgfarbe\,#2}}}}\luecke[\lsgbreite]\else\luecke[#1]\fi}
% Schreiblinien mit Lösung: \lsglinien{Anzahl}{Lösungstext} (gleiche Höhe wie \linien:
%  je Linie \baselineskip + 5,5 mm + \parskip, ausgemessen)
\newlength{\lsglinienhoehe}
\newcommand{\lsglinien}[2]{\ifloesung\par\vspace{1.5mm}\setlength{\lsglinienhoehe}{\dimexpr #1\dimexpr\baselineskip+5.5mm+\parskip\relax-16.25pt\relax}\noindent\begin{minipage}[t][\lsglinienhoehe][t]{\linewidth}\lsgfarbe #2\end{minipage}\par\vspace{-1.5mm}\else\linien{#1}\fi}
% Beliebiger Inhalt: \lsgoder{nur in der Lösung}{nur im Blatt}
\newcommand{\lsgoder}[2]{\ifloesung\leavevmode#1\else#2\fi}
% Text nur in der Lösung sichtbar, belegt im Blatt aber denselben Platz (ohne Text im PDF)
\newcommand{\lsgtext}[1]{\leavevmode\ifloesung{\lsgfarbe#1}\else\phantom{#1}\fi}
% Ankreuzfeld: \kreuz{1} = angekreuzt (nur in der Lösung), \kreuz{0} = leer
\newcommand{\kreuz}[1]{\leavevmode\ifloesung\ifnum#1=1 {\lsgfarbe$\boxtimes$}\else$\square$\fi\else$\square$\fi}
% Tabellenzeile mit fester Höhe für Lösungseinträge (2 Zeilen \small)
\newcommand{\zeile}{\rule[-6mm]{0pt}{9.5mm}}

% ---------- Verkettete Strukturen zeichnen ----------
% Objekt mit Kopf und Attributen:  \node[obj] (k1) {k1: Knoten \nodepart{second} daten = Hans\\ nachfolger};
\tikzset{
  obj/.style={draw,rectangle split,rectangle split parts=2,rectangle split part fill={orange!25,orange!8},
              rounded corners=2pt,font=\scriptsize,inner sep=2pt,align=left},
  qobj/.style={obj,rectangle split part fill={teal!25,teal!8}},
  aobj/.style={obj,rectangle split part fill={gray!30,gray!8}},
  pobj/.style={draw,rounded corners=2pt,fill=gray!8,font=\scriptsize,inner sep=2pt,align=left},
  klasse/.style={draw,rectangle split,rectangle split parts=3,font=\scriptsize,align=left,fill=white},
  zeiger/.style={thick,-{Stealth[length=2.2mm]}},
  erbt/.style={thick,-{Triangle[open,length=2.6mm,width=2.4mm]}},
  implementiert/.style={thick,densely dashed,-{Triangle[open,length=2.6mm,width=2.4mm]}},
  lsgzeiger/.style={zeiger,red!75!black},
  nullref/.style={font=\scriptsize\ttfamily,inner sep=1pt},
  schritt/.style={circle,draw,fill=white,inner sep=0.6pt,font=\scriptsize\bfseries},
  lsgschritt/.style={schritt,draw=red!75!black,text=red!75!black},
  % Sequenzdiagramm (UML 2): Kopf eckig, synchron = gefüllte Spitze, asynchron = offene Spitze,
  %  Rückgabe und Erzeugung gestrichelt mit offener Spitze (kein hohles Dreieck: das wäre Vererbung)
  lebenslinie/.style={draw,sharp corners,fill=white,font=\scriptsize,inner sep=2pt},
  leben/.style={densely dashed,gray},
  aufruf/.style={thick,-{Stealth[length=2mm]}},
  asynchron/.style={thick,-{Straight Barb[length=2mm,width=2.4mm]}},
  rueckgabe/.style={thick,densely dashed,-{Straight Barb[length=2mm,width=2.4mm]}},
  erzeugung/.style={rueckgabe},
}
% Akteur (Strichmännchen) mit Namen darunter: \akteur{x}{y oben}{Name}; ca. 0,9 cm hoch
\newcommand{\akteur}[3]{\begin{scope}[shift={(#1,#2)},line width=0.6pt]
  \draw (0,-0.09cm) circle (0.09cm);
  \draw (0,-0.18cm) -- (0,-0.42cm) (-0.16cm,-0.26cm) -- (0.16cm,-0.26cm)
        (-0.13cm,-0.62cm) -- (0,-0.42cm) -- (0.13cm,-0.62cm);
  \node[font=\scriptsize,inner sep=1pt,anchor=north] at (0,-0.63cm) {#3};
\end{scope}}
% Aktivierungsbalken auf der Lebenslinie bei x von y1 bis y2: \aktiv{x}{y1}{y2}
\newcommand{\aktiv}[3]{\draw[fill=gray!15] ($(#1,#2)+(-0.09cm,0)$) rectangle ($(#1,#3)+(0.09cm,0)$);}
% Lösungszeichnung: im Blatt unsichtbar, belegt aber denselben Platz
\ifloesung
  \tikzset{lsgbild/.style={red!75!black}}
\else
  \tikzset{lsgbild/.style={draw opacity=0,fill opacity=0,text opacity=0}}
\fi

% ---------- Java-Listings ----------
\lstset{language=Java,basicstyle=\ttfamily\small,keywordstyle=\bfseries,commentstyle=\color{gray!80!black},stringstyle=\color{black},showstringspaces=false,columns=flexible,keepspaces=true,tabsize=3,frame=single,framerule=0.4pt,rulecolor=\color{black},xleftmargin=2mm,framexleftmargin=1mm,aboveskip=2mm,belowskip=2mm,breaklines=true,escapechar=§,
 morekeywords={interface,implements,abstract,extends},
 literate={ä}{{\"a}}1 {ö}{{\"o}}1 {ü}{{\"u}}1 {Ä}{{\"A}}1 {Ö}{{\"O}}1 {Ü}{{\"U}}1 {ß}{{\ss}}1 {–}{{--}}1 {„}{{\glqq}}1 {“}{{\grqq}}1 {→}{{$\to$}}1}
\lstnewenvironment{java}[1][]{\lstset{#1}}{}

% ---------- Antwortplatz nach CLAUDE.md (in Kapitel 2 immer an) ----------
\def\antwortplatzacht{}
% Makros wie informatik/10/java/arbeitsblaetter/ab-vorlage.tex:
%  \linien{n}/\lsglinien{n}{Text}  Schreiblinien 8 mm      \linienfeld{n}/\lsgfeld{n}{Code}
%  \luecke[l]/\lsg[l]{Text}         Lücke (Standard 6,75 cm, Fachbegriff 5,25 cm, Zahl 2,25 cm)
%  \restzeile/\lsgrest{Text}        Lücke bis Zeilenende    \lueckentext  Absatz mit Lücken
%  \schreibzeile                    Tabellenzelle 8 mm      \lsgtext{…}  nur in der Lösung
%  \kariert{n}/\lsgkariert{n}{TikZ}, \karobox{b}{h}/\lsgkarobox{b}{h}{TikZ}  5-mm-Karo, ganze Kästchen
%  javaluecke                       Listing mit 1,5-fachem Zeilenabstand für Lücken (§\luecke§)
\ifdefined\antwortplatzacht
\usepackage{xspace}
\newlength{\zeilenabstand}\setlength{\zeilenabstand}{8mm}
\newlength{\karo}\setlength{\karo}{5mm}
\tikzset{karo/.style={draw=gray!50,line width=0.3pt},schreiblinie/.style={draw=black!65,line width=0.35pt}}
\renewcommand{\lsg}[2][6.75cm]{\leavevmode\ifloesung\settowidth{\lsgbreite}{\,#2\,}%
  \ifdim\lsgbreite>#1\relax\rlap{\raisebox{0.8pt}{\resizebox{#1}{!}{\lsgfarbe\,#2\,}}}%
  \else\rlap{\raisebox{0.8pt}{\lsgfarbe\,#2}}\fi\fi\rule[-0.6ex]{#1}{0.4pt}\xspace}
\renewcommand{\luecke}[1][6.75cm]{\lsg[#1]{}}
\newcommand{\lsgrest}[1]{\leavevmode\unskip\hspace{1.5mm}\ifloesung\rlap{\raisebox{0.8pt}{\lsgfarbe\,#1}}\fi\leaders\hbox{\rule[-0.6ex]{0.5pt}{0.4pt}}\hfill\null\par}
\newcommand{\restzeile}{\lsgrest{}}
\renewcommand{\lsgtext}[1]{\ifloesung\leavevmode{\lsgfarbe#1}\fi}
\newcommand{\lueckentext}{\linespread{1.5}\selectfont\raggedright}
\newcommand{\schreibzeile}{\rule[-2.8mm]{0pt}{8mm}}
\renewcommand{\lsglinien}[2]{\par\vspace{0.5mm}\noindent\begin{tikzpicture}
  \path (0,0) (\linewidth-0.4pt,-#1\zeilenabstand-1.5mm);
  \foreach \i in {1,...,#1}{\draw[schreiblinie] (0.2pt,-\i\zeilenabstand) -- (\linewidth-0.2pt,-\i\zeilenabstand);}
  \ifloesung\node[overlay,anchor=north west,inner sep=0pt] at (1mm,-\zeilenabstand+1pt+5mm) {\parbox[t]{\linewidth-2mm}{\lsgfarbe\raggedright\setlength{\baselineskip}{\zeilenabstand}\rule{0pt}{5mm}#2}};\fi
  \end{tikzpicture}\par}
\renewcommand{\linien}[1]{\lsglinien{#1}{}}
\newcommand{\lsgfeld}[2]{\noindent\begin{tikzpicture}
  \draw[line width=0.4pt] (0.2pt,0) rectangle (\linewidth-0.2pt,-#1\zeilenabstand-3mm);
  \foreach \i in {1,...,#1}{\draw[schreiblinie] (2mm,-\i\zeilenabstand) -- (\linewidth-2mm,-\i\zeilenabstand);}
  \ifloesung\node[overlay,anchor=north west,inner sep=0pt] at (2.5mm,-\zeilenabstand+1pt+5mm) {\parbox[t]{\linewidth-5mm}{\lsgfarbe\ttfamily\small\setlength{\baselineskip}{\zeilenabstand}\rule{0pt}{5mm}#2}};\fi
  \end{tikzpicture}}
\newcommand{\linienfeld}[1]{\lsgfeld{#1}{}}
\newcommand{\karoraster}[3]{\begin{tikzpicture}\draw[karo] (0,0) grid[step=\karo] (#1\karo,#2\karo);
  \ifloesung\begin{scope}[overlay,shift={(0,#2\karo)},red!75!black,text=red!75!black,line width=0.8pt]#3\end{scope}\fi\end{tikzpicture}}
\newcommand{\lsgkariert}[2]{\par\vspace{1.5mm}\noindent\pgfmathtruncatemacro{\karoX}{floor(\linewidth/\karo)}%
  \makebox[\linewidth]{\karoraster{\karoX}{#1}{#2}}\par\vspace{1mm}}
\renewcommand{\kariert}[1]{\lsgkariert{#1}{}}
\newcommand{\lsgkarobox}[3]{\pgfmathtruncatemacro{\karoX}{round((#1)/\karo)}\pgfmathtruncatemacro{\karoY}{round((#2)/\karo)}\karoraster{\karoX}{\karoY}{#3}}
\renewcommand{\karobox}[2]{\lsgkarobox{#1}{#2}{}}
\lstnewenvironment{javaluecke}[1][]{\lstset{basicstyle=\linespread{1.5}\small\ttfamily,#1}}{}
\fi

% ---------- Binärbäume ----------
% Knoten (bk), Abschluss (ab), markierter Knoten (bkm), Lösungsknoten (lsgbk, im Blatt unsichtbar)
\tikzset{
  baum/.style={level distance=9mm,level 1/.style={sibling distance=36mm},level 2/.style={sibling distance=18mm},
               level 3/.style={sibling distance=9mm},level 4/.style={sibling distance=6mm},
               edge from parent/.style={draw,thick,gray!60!black}},
  bk/.style={circle,draw=orange!70!black,fill=orange!22,minimum size=7.6mm,inner sep=0pt,font=\scriptsize,thick},
  bkm/.style={bk,fill=yellow!40,draw=orange!90!black,line width=1.2pt},
  ab/.style={rectangle,draw=gray!70,fill=gray!25,minimum size=2.2mm,inner sep=0pt},
  hoehe/.style={font=\scriptsize\bfseries,inner sep=1pt},
}
% Das Spielerverzeichnis (Mika, Emil, Tara, Bo, Ines, Pia, Zoe, Dana, Kai, Ruby, Vik); Argument: TikZ-Optionen
\newcommand{\verzeichnisbaum}[1][]{\begin{tikzpicture}[baum,#1]
\node[bk] (Mika) {Mika}
  child {node[bk] (Emil) {Emil}
    child {node[bk] (Bo) {Bo} child[missing] child {node[bk] (Dana) {Dana}}}
    child {node[bk] (Ines) {Ines} child[missing] child {node[bk] (Kai) {Kai}}}}
  child {node[bk] (Tara) {Tara}
    child {node[bk] (Pia) {Pia} child[missing] child {node[bk] (Ruby) {Ruby}}}
    child {node[bk] (Zoe) {Zoe} child {node[bk] (Vik) {Vik}} child[missing]}};
\end{tikzpicture}}

\begin{document}
\abtitel{2.1}{Bäume und geordnete Binärbäume}

Das Online-Spiel \textbf{Sternenstaub} hat zwei Millionen Accounts. Beim Login sucht der Server den Namen bisher in einer Liste. Du baust ein schnelleres Spielerverzeichnis. \web{Station 2.1}

\begin{aufgabe}[Namen raten]{1}
\begin{lueckentext}
\textbf{a)} Ich habe \lsg[2.25cm]{z.\,B. 8} Fragen gebraucht. Meine Strategie:\lsgrest{immer nach dem Namen in der Mitte fragen.}
\textbf{b)} Jede Frage nach der Mitte \lsg[5.25cm]{halbiert} den Bereich. Für 1000 Namen reichen deshalb höchstens \lsg[2.25cm]{10} Fragen, denn $2^{10} =$ \lsg[2.25cm]{1024}. Bei 2 Millionen Accounts sind es höchstens \lsg[2.25cm]{21} Fragen, eine Liste braucht bis zu \lsg[3.6cm]{2\,000\,000} Vergleiche.
\end{lueckentext}
\end{aufgabe}

\begin{merke}[Merke: Bäume]
\begin{lueckentext}
Ein Baum besteht aus \lsg[5.25cm]{Knoten}, die durch \lsg[5.25cm]{Kanten} verbunden sind. Der oberste Knoten heißt \lsg[5.25cm]{Wurzel}. Jeder andere Knoten hat genau einen \lsg[5.25cm]{Elternknoten}. Knoten ohne Kinder heißen \lsg[5.25cm]{Blätter}, alle anderen \lsg[5.25cm]{innere Knoten}. In einem \textbf{Binärbaum} hat jeder Knoten höchstens \lsg[2.25cm]{zwei} Kinder. Jeder Knoten bildet mit allem, was unter ihm hängt, einen \lsg[5.25cm]{Teilbaum}. Die Wurzel liegt auf Ebene 1. Die Anzahl der Ebenen ist die \lsg[5.25cm]{Höhe} des Baums.\par
\textbf{Rekursiv:} Ein Binärbaum ist entweder \lsg[5.25cm]{leer} oder er besteht aus einer Wurzel mit einem linken und einem rechten Teilbaum, die selbst wieder \lsg[5.25cm]{Binärbäume} sind.
\end{lueckentext}
\end{merke}

\begin{aufgabe}[Baum-Safari: das Spielerverzeichnis]{2}
\begin{minipage}[t]{0.5\linewidth}
\vspace{0pt}\centering
\verzeichnisbaum[scale=0.98,every node/.append style={transform shape}]
\end{minipage}\hfill
\begin{minipage}[t]{0.47\linewidth}
\vspace{0pt}\lueckentext
Wurzel: \lsg[3.6cm]{Mika}\par
Blätter: \lsgrest{Dana, Kai, Ruby, Vik}
Elternknoten von Kai: \lsg[2.6cm]{Ines}\par
Kinder von Tara: \lsg[4.2cm]{Pia, Zoe}\par
Rechter Teilbaum von Emil: \lsgrest{Ines, Kai}
Knoten auf Ebene 3: \lsgrest{Bo, Ines, Pia, Zoe}
Höhe des Baums: \lsg[2.25cm]{4}
\end{minipage}
\end{aufgabe}

\begin{merke}[Merke: Geordneter Binärbaum (Suchbaum)]
\begin{lueckentext}
Für \textbf{jeden} Knoten gilt: Alle Schlüssel im linken Teilbaum sind \lsg[5.25cm]{kleiner}, alle im rechten Teilbaum sind \lsg[5.25cm]{größer} als sein eigener Schlüssel. Jeder Schlüssel kommt nur einmal vor. (Hier: Schlüssel = Name, „kleiner“ = kommt im Alphabet früher.)\par
\end{lueckentext}
\begin{minipage}[t]{0.49\linewidth}
\textbf{Vorgehensweise Suchen} (ab der Wurzel):\\
\nr{1} gleich: gefunden.\\
\nr{2} kleiner: im linken Teilbaum weitersuchen.\\
\nr{3} größer: im rechten Teilbaum weitersuchen.\\
\nr{4} Teilbaum leer: Den Namen gibt es nicht.
\end{minipage}\hfill
\begin{minipage}[t]{0.49\linewidth}
\textbf{Vorgehensweise Einfügen} (ab der Wurzel):\\
\nr{1} gleich: Name vergeben, nicht einfügen.\\
\nr{2} kleiner: im linken Teilbaum einfügen.\\
\nr{3} größer: im rechten Teilbaum einfügen.\\
\nr{4} Teilbaum leer: Hier entsteht der neue Knoten.
\end{minipage}\par\vspace{1mm}
\begin{lueckentext}
Der neue Knoten ist immer ein \lsg[5.25cm]{Blatt}. Die Suche geht bei jedem Vergleich in nur \lsg[2.25cm]{einen} Teilbaum, sie braucht also höchstens so viele Vergleiche, wie der Baum \lsg[5.25cm]{Ebenen} hat.
\end{lueckentext}
\end{merke}

\newpage
\begin{aufgabe}[Such-Sprint: Suchwege im Verzeichnis (Baum aus Aufgabe 2)]{3}
{\small\renewcommand{\arraystretch}{1.1}
\begin{tabularx}{\linewidth}{|l|X|c|c|}\hline
gesucht & Suchweg (besuchte Knoten) & Vergleiche & gefunden? \\\hline
Ruby & \schreibzeile\lsgtext{Mika → Tara → Pia → Ruby} & \lsgtext{4} & \lsgtext{ja} \\\hline
Dana & \schreibzeile\lsgtext{Mika → Emil → Bo → Dana} & \lsgtext{4} & \lsgtext{ja} \\\hline
Olaf & \schreibzeile\lsgtext{Mika → Tara → Pia → (links unter Pia leer)} & \lsgtext{3} & \lsgtext{nein} \\\hline
\end{tabularx}}
\end{aufgabe}

\begin{aufgabe}[Baumschule]{4}
\textbf{a)} Zeichne den Baum zur Reihenfolge \code{Ole Kim Tom Ben Lia Sam Zara}.\hfill Höhe: \lsg[2.25cm]{3}\par\vspace{1mm}
\lsgkariert{6}{
  \begin{scope}[baum,level distance=9mm,level 1/.style={sibling distance=44mm},level 2/.style={sibling distance=20mm},every node/.append style={draw=red!75!black,fill=white}]
  \node[bk] at (8.5,-0.6) {Ole}
    child {node[bk] {Kim} child {node[bk] {Ben}} child {node[bk] {Lia}}}
    child {node[bk] {Tom} child {node[bk] {Sam}} child {node[bk] {Zara}}};
  \end{scope}}\par\vspace{1mm}
\textbf{b)} Skizziere den Baum für dieselben Namen alphabetisch.\hfill Höhe: \lsg[2.25cm]{7}\par\vspace{1mm}
\begin{minipage}[t]{0.52\linewidth}
\vspace{0pt}
\lsgkarobox{8.5cm}{4cm}{
  \foreach \n/\i in {Ben/0,Kim/1,Lia/2,Ole/3,Sam/4,Tom/5,Zara/6}{
    \node[bk,minimum size=6mm,font=\tiny,draw=red!75!black,fill=white] (z\i) at (0.6+1.2*\i,-0.45-0.5*\i) {\n};}
  \foreach \i [evaluate=\i as \j using int(\i+1)] in {0,...,5}{\draw[thick] (z\i) -- (z\j);}}
\end{minipage}\hfill
\begin{minipage}[t]{0.45\linewidth}
\vspace{0pt}\lueckentext
Werden die Namen sortiert eingefügt,\par\lsgrest{hängt jeder neue Name ganz rechts:}
\lsgrest{eine Kette, so langsam wie eine Liste}
\lsgrest{(entarteter Baum).}
\end{minipage}
\end{aufgabe}

\begin{aufgabe}[Ist das ein geordneter Binärbaum?]{5}
\begin{minipage}[t]{0.38\linewidth}
\vspace{0pt}\centering
\begin{tikzpicture}[baum,level distance=8mm,level 1/.style={sibling distance=30mm},level 2/.style={sibling distance=15mm},level 3/.style={sibling distance=8mm},every node/.append style={scale=0.9}]
\node[bk] {Mika}
  child {node[bk] {Emil} child {node[bk] {Bo}} child {node[bk] {Ines}}}
  child {node[bk] {Tara} child {node[bk] {Pia} child {node[bk] {Kai}} child[missing]} child {node[bk] {Zoe}}};
\end{tikzpicture}
\end{minipage}\hfill
\begin{minipage}[t]{0.6\linewidth}
\vspace{0pt}
Jemand hat Kai von Hand links unter Pia eingetragen, „weil Kai vor Pia kommt“.\par\vspace{1mm}
\kreuz{0} ja, geordnet\qquad \kreuz{1} nein, nicht geordnet, weil:\par
\lsglinien{3}{Kai steht im rechten Teilbaum von Mika, kommt aber vor Mika. Die Suche nach Kai biegt bei Mika links ab und findet ihn nie.}
\vspace{2mm}\begin{lueckentext}
Kai gehört \lsgrest{rechts unter Ines.}
\end{lueckentext}
\end{minipage}
\end{aufgabe}
\end{document}
;
%            füllt alle Lücken rot aus; siehe \lsg, \lsglinien, \lsgoder,
%            \lsgtext, lsgbild, \kreuz)
% =====================================================================
\usepackage[a4paper,left=16mm,right=16mm,top=23mm,bottom=13mm,headheight=12mm,headsep=3mm,footskip=6mm]{geometry}
\usepackage{iftex}
\ifPDFTeX
  \usepackage[T1]{fontenc}
  \usepackage[utf8]{inputenc}
  \usepackage{helvet}
  \usepackage[scaled=0.86]{beramono}
  \renewcommand{\familydefault}{\sfdefault}
\else
  \usepackage{fontspec}
  \setmainfont{texgyreheros}[Extension=.otf,UprightFont=*-regular,BoldFont=*-bold,ItalicFont=*-italic,BoldItalicFont=*-bolditalic]
  \setmonofont{DejaVuSansMono}[Extension=.ttf,UprightFont=*,BoldFont=*-Bold,ItalicFont=*-Oblique,BoldItalicFont=*-BoldOblique,Scale=0.86]
\fi
\usepackage[ngerman,shorthands=off]{babel}
\usepackage{amssymb}
\usepackage{xcolor,graphicx,tabularx,array,enumitem,fancyhdr,tikz,listings,multicol,calc,etoolbox,pifont}
\usepackage[most]{tcolorbox}
\usetikzlibrary{arrows.meta,positioning,calc,decorations.pathreplacing,shapes.multipart}
\usepackage[hidelinks]{hyperref}
\setlength{\parindent}{0pt}
\setlength{\parskip}{3pt}
\setlist{nosep,leftmargin=*}

% ---------- Lösungsschalter ----------
\newif\ifloesung
\ifdefined\mitloesung\loesungtrue\fi
\newcommand{\lsgfarbe}{\color{red!75!black}}

% ---------- Kopf- und Fußzeile (einheitliche Form) ----------
\newcommand{\lehrkraft}{Hau}
\newcommand{\themenbereich}{2. Binärbäume und Tiefensuche}
\newcommand{\abnummer}{}
\pagestyle{fancy}\fancyhf{}
\renewcommand{\headrulewidth}{0.4pt}
\fancyhead[L]{\small Klasse 12\\Datum:}
\fancyhead[C]{\small Themenbereich:\\\themenbereich}
\fancyhead[R]{\small \lehrkraft\ifloesung\\{\lsgfarbe\bfseries Lösung}\fi}
\fancyfoot[C]{\scriptsize Informatik 12\,·\,Blatt \abnummer\ifloesung\ (Lösung)\fi\,·\,Seite \thepage}
\newcommand{\abtitel}[2]{\renewcommand{\abnummer}{#1}\begin{center}\Large\bfseries #1\quad\underline{#2}\end{center}\vspace{-2mm}}

% ---------- Boxen ----------
\newenvironment{aufgabe}[2][]{\begin{tcolorbox}[enhanced,breakable,colback=white,colframe=black,boxrule=0.7pt,arc=3mm,left=2mm,right=2mm,top=1mm,bottom=1.2mm,before skip=2mm,after skip=2mm]\textbf{Aufgabe #2}\ifstrempty{#1}{}{ \textit{#1}}\par\vspace{0.3mm}}{\end{tcolorbox}}
\newtcolorbox{merke}[1][Merke]{enhanced,breakable,colback=green!5,colframe=green!40!black,boxrule=0.8pt,arc=2mm,left=2mm,right=2mm,top=1mm,bottom=1.2mm,before skip=2mm,after skip=2mm,title={#1},fonttitle=\bfseries,coltitle=white,toptitle=0.3mm,bottomtitle=0.3mm}
\newtcolorbox{hinweis}{enhanced,breakable,colback=yellow!12,colframe=orange!70!black,boxrule=0.6pt,arc=2mm,left=2mm,right=2mm,top=1mm,bottom=1mm,before skip=2mm,after skip=2mm}
\newcommand{\web}[1]{{\small\color{gray}\textit{Website: #1}}}

% ---------- Lücken, Linien, Karos ----------
\newcommand{\luecke}[1][3.5cm]{\rule[-0.5ex]{#1}{0.4pt}}
\newcommand{\linien}[1]{\par\vspace{1.5mm}\foreach \i in {1,...,#1}{\noindent\rule[0pt]{\linewidth}{0.3pt}\par\vspace{5.5mm}}\vspace{-3mm}}
\newcommand{\kariert}[1]{\par\vspace{1.5mm}\noindent\begin{tikzpicture}\draw[step=5mm,gray!55,thin] (0,0) grid (\linewidth-1pt,#1*5mm);\end{tikzpicture}\par}
\newcommand{\karobox}[2]{\begin{tikzpicture}\draw[step=5mm,gray!55,thin] (0,0) grid (#1,#2);\end{tikzpicture}}
\newcommand{\code}[1]{\texttt{#1}}
% Schritt-Nummern wie im Code und auf der Website: \nr{1} = ①
\newcommand{\nr}[1]{\ding{\numexpr171+#1\relax}}

% Lücke mit Lösung: \lsg[Mindestbreite]{Lösungstext}
%  Die Lösung steht (ohne eigene Höhe/Tiefe) auf derselben Linie wie im Blatt: Das Layout verschiebt sich nicht,
%  solange der Text nicht breiter als die Mindestbreite ist.
\newlength{\lsgbreite}
\newcommand{\lsg}[2][3.5cm]{\leavevmode\ifloesung\settowidth{\lsgbreite}{\,#2\,}\ifdim\lsgbreite<#1\setlength{\lsgbreite}{#1}\fi\rlap{\smash{\raisebox{0.8pt}{\makebox[\lsgbreite][l]{\lsgfarbe\,#2}}}}\luecke[\lsgbreite]\else\luecke[#1]\fi}
% Schreiblinien mit Lösung: \lsglinien{Anzahl}{Lösungstext} (gleiche Höhe wie \linien:
%  je Linie \baselineskip + 5,5 mm + \parskip, ausgemessen)
\newlength{\lsglinienhoehe}
\newcommand{\lsglinien}[2]{\ifloesung\par\vspace{1.5mm}\setlength{\lsglinienhoehe}{\dimexpr #1\dimexpr\baselineskip+5.5mm+\parskip\relax-16.25pt\relax}\noindent\begin{minipage}[t][\lsglinienhoehe][t]{\linewidth}\lsgfarbe #2\end{minipage}\par\vspace{-1.5mm}\else\linien{#1}\fi}
% Beliebiger Inhalt: \lsgoder{nur in der Lösung}{nur im Blatt}
\newcommand{\lsgoder}[2]{\ifloesung\leavevmode#1\else#2\fi}
% Text nur in der Lösung sichtbar, belegt im Blatt aber denselben Platz (ohne Text im PDF)
\newcommand{\lsgtext}[1]{\leavevmode\ifloesung{\lsgfarbe#1}\else\phantom{#1}\fi}
% Ankreuzfeld: \kreuz{1} = angekreuzt (nur in der Lösung), \kreuz{0} = leer
\newcommand{\kreuz}[1]{\leavevmode\ifloesung\ifnum#1=1 {\lsgfarbe$\boxtimes$}\else$\square$\fi\else$\square$\fi}
% Tabellenzeile mit fester Höhe für Lösungseinträge (2 Zeilen \small)
\newcommand{\zeile}{\rule[-6mm]{0pt}{9.5mm}}

% ---------- Verkettete Strukturen zeichnen ----------
% Objekt mit Kopf und Attributen:  \node[obj] (k1) {k1: Knoten \nodepart{second} daten = Hans\\ nachfolger};
\tikzset{
  obj/.style={draw,rectangle split,rectangle split parts=2,rectangle split part fill={orange!25,orange!8},
              rounded corners=2pt,font=\scriptsize,inner sep=2pt,align=left},
  qobj/.style={obj,rectangle split part fill={teal!25,teal!8}},
  aobj/.style={obj,rectangle split part fill={gray!30,gray!8}},
  pobj/.style={draw,rounded corners=2pt,fill=gray!8,font=\scriptsize,inner sep=2pt,align=left},
  klasse/.style={draw,rectangle split,rectangle split parts=3,font=\scriptsize,align=left,fill=white},
  zeiger/.style={thick,-{Stealth[length=2.2mm]}},
  erbt/.style={thick,-{Triangle[open,length=2.6mm,width=2.4mm]}},
  implementiert/.style={thick,densely dashed,-{Triangle[open,length=2.6mm,width=2.4mm]}},
  lsgzeiger/.style={zeiger,red!75!black},
  nullref/.style={font=\scriptsize\ttfamily,inner sep=1pt},
  schritt/.style={circle,draw,fill=white,inner sep=0.6pt,font=\scriptsize\bfseries},
  lsgschritt/.style={schritt,draw=red!75!black,text=red!75!black},
  % Sequenzdiagramm (UML 2): Kopf eckig, synchron = gefüllte Spitze, asynchron = offene Spitze,
  %  Rückgabe und Erzeugung gestrichelt mit offener Spitze (kein hohles Dreieck: das wäre Vererbung)
  lebenslinie/.style={draw,sharp corners,fill=white,font=\scriptsize,inner sep=2pt},
  leben/.style={densely dashed,gray},
  aufruf/.style={thick,-{Stealth[length=2mm]}},
  asynchron/.style={thick,-{Straight Barb[length=2mm,width=2.4mm]}},
  rueckgabe/.style={thick,densely dashed,-{Straight Barb[length=2mm,width=2.4mm]}},
  erzeugung/.style={rueckgabe},
}
% Akteur (Strichmännchen) mit Namen darunter: \akteur{x}{y oben}{Name}; ca. 0,9 cm hoch
\newcommand{\akteur}[3]{\begin{scope}[shift={(#1,#2)},line width=0.6pt]
  \draw (0,-0.09cm) circle (0.09cm);
  \draw (0,-0.18cm) -- (0,-0.42cm) (-0.16cm,-0.26cm) -- (0.16cm,-0.26cm)
        (-0.13cm,-0.62cm) -- (0,-0.42cm) -- (0.13cm,-0.62cm);
  \node[font=\scriptsize,inner sep=1pt,anchor=north] at (0,-0.63cm) {#3};
\end{scope}}
% Aktivierungsbalken auf der Lebenslinie bei x von y1 bis y2: \aktiv{x}{y1}{y2}
\newcommand{\aktiv}[3]{\draw[fill=gray!15] ($(#1,#2)+(-0.09cm,0)$) rectangle ($(#1,#3)+(0.09cm,0)$);}
% Lösungszeichnung: im Blatt unsichtbar, belegt aber denselben Platz
\ifloesung
  \tikzset{lsgbild/.style={red!75!black}}
\else
  \tikzset{lsgbild/.style={draw opacity=0,fill opacity=0,text opacity=0}}
\fi

% ---------- Java-Listings ----------
\lstset{language=Java,basicstyle=\ttfamily\small,keywordstyle=\bfseries,commentstyle=\color{gray!80!black},stringstyle=\color{black},showstringspaces=false,columns=flexible,keepspaces=true,tabsize=3,frame=single,framerule=0.4pt,rulecolor=\color{black},xleftmargin=2mm,framexleftmargin=1mm,aboveskip=2mm,belowskip=2mm,breaklines=true,escapechar=§,
 morekeywords={interface,implements,abstract,extends},
 literate={ä}{{\"a}}1 {ö}{{\"o}}1 {ü}{{\"u}}1 {Ä}{{\"A}}1 {Ö}{{\"O}}1 {Ü}{{\"U}}1 {ß}{{\ss}}1 {–}{{--}}1 {„}{{\glqq}}1 {“}{{\grqq}}1 {→}{{$\to$}}1}
\lstnewenvironment{java}[1][]{\lstset{#1}}{}

% ---------- Antwortplatz nach CLAUDE.md (in Kapitel 2 immer an) ----------
\def\antwortplatzacht{}
% Makros wie informatik/10/java/arbeitsblaetter/ab-vorlage.tex:
%  \linien{n}/\lsglinien{n}{Text}  Schreiblinien 8 mm      \linienfeld{n}/\lsgfeld{n}{Code}
%  \luecke[l]/\lsg[l]{Text}         Lücke (Standard 6,75 cm, Fachbegriff 5,25 cm, Zahl 2,25 cm)
%  \restzeile/\lsgrest{Text}        Lücke bis Zeilenende    \lueckentext  Absatz mit Lücken
%  \schreibzeile                    Tabellenzelle 8 mm      \lsgtext{…}  nur in der Lösung
%  \kariert{n}/\lsgkariert{n}{TikZ}, \karobox{b}{h}/\lsgkarobox{b}{h}{TikZ}  5-mm-Karo, ganze Kästchen
%  javaluecke                       Listing mit 1,5-fachem Zeilenabstand für Lücken (§\luecke§)
\ifdefined\antwortplatzacht
\usepackage{xspace}
\newlength{\zeilenabstand}\setlength{\zeilenabstand}{8mm}
\newlength{\karo}\setlength{\karo}{5mm}
\tikzset{karo/.style={draw=gray!50,line width=0.3pt},schreiblinie/.style={draw=black!65,line width=0.35pt}}
\renewcommand{\lsg}[2][6.75cm]{\leavevmode\ifloesung\settowidth{\lsgbreite}{\,#2\,}%
  \ifdim\lsgbreite>#1\relax\rlap{\raisebox{0.8pt}{\resizebox{#1}{!}{\lsgfarbe\,#2\,}}}%
  \else\rlap{\raisebox{0.8pt}{\lsgfarbe\,#2}}\fi\fi\rule[-0.6ex]{#1}{0.4pt}\xspace}
\renewcommand{\luecke}[1][6.75cm]{\lsg[#1]{}}
\newcommand{\lsgrest}[1]{\leavevmode\unskip\hspace{1.5mm}\ifloesung\rlap{\raisebox{0.8pt}{\lsgfarbe\,#1}}\fi\leaders\hbox{\rule[-0.6ex]{0.5pt}{0.4pt}}\hfill\null\par}
\newcommand{\restzeile}{\lsgrest{}}
\renewcommand{\lsgtext}[1]{\ifloesung\leavevmode{\lsgfarbe#1}\fi}
\newcommand{\lueckentext}{\linespread{1.5}\selectfont\raggedright}
\newcommand{\schreibzeile}{\rule[-2.8mm]{0pt}{8mm}}
\renewcommand{\lsglinien}[2]{\par\vspace{0.5mm}\noindent\begin{tikzpicture}
  \path (0,0) (\linewidth-0.4pt,-#1\zeilenabstand-1.5mm);
  \foreach \i in {1,...,#1}{\draw[schreiblinie] (0.2pt,-\i\zeilenabstand) -- (\linewidth-0.2pt,-\i\zeilenabstand);}
  \ifloesung\node[overlay,anchor=north west,inner sep=0pt] at (1mm,-\zeilenabstand+1pt+5mm) {\parbox[t]{\linewidth-2mm}{\lsgfarbe\raggedright\setlength{\baselineskip}{\zeilenabstand}\rule{0pt}{5mm}#2}};\fi
  \end{tikzpicture}\par}
\renewcommand{\linien}[1]{\lsglinien{#1}{}}
\newcommand{\lsgfeld}[2]{\noindent\begin{tikzpicture}
  \draw[line width=0.4pt] (0.2pt,0) rectangle (\linewidth-0.2pt,-#1\zeilenabstand-3mm);
  \foreach \i in {1,...,#1}{\draw[schreiblinie] (2mm,-\i\zeilenabstand) -- (\linewidth-2mm,-\i\zeilenabstand);}
  \ifloesung\node[overlay,anchor=north west,inner sep=0pt] at (2.5mm,-\zeilenabstand+1pt+5mm) {\parbox[t]{\linewidth-5mm}{\lsgfarbe\ttfamily\small\setlength{\baselineskip}{\zeilenabstand}\rule{0pt}{5mm}#2}};\fi
  \end{tikzpicture}}
\newcommand{\linienfeld}[1]{\lsgfeld{#1}{}}
\newcommand{\karoraster}[3]{\begin{tikzpicture}\draw[karo] (0,0) grid[step=\karo] (#1\karo,#2\karo);
  \ifloesung\begin{scope}[overlay,shift={(0,#2\karo)},red!75!black,text=red!75!black,line width=0.8pt]#3\end{scope}\fi\end{tikzpicture}}
\newcommand{\lsgkariert}[2]{\par\vspace{1.5mm}\noindent\pgfmathtruncatemacro{\karoX}{floor(\linewidth/\karo)}%
  \makebox[\linewidth]{\karoraster{\karoX}{#1}{#2}}\par\vspace{1mm}}
\renewcommand{\kariert}[1]{\lsgkariert{#1}{}}
\newcommand{\lsgkarobox}[3]{\pgfmathtruncatemacro{\karoX}{round((#1)/\karo)}\pgfmathtruncatemacro{\karoY}{round((#2)/\karo)}\karoraster{\karoX}{\karoY}{#3}}
\renewcommand{\karobox}[2]{\lsgkarobox{#1}{#2}{}}
\lstnewenvironment{javaluecke}[1][]{\lstset{basicstyle=\linespread{1.5}\small\ttfamily,#1}}{}
\fi

% ---------- Binärbäume ----------
% Knoten (bk), Abschluss (ab), markierter Knoten (bkm), Lösungsknoten (lsgbk, im Blatt unsichtbar)
\tikzset{
  baum/.style={level distance=9mm,level 1/.style={sibling distance=36mm},level 2/.style={sibling distance=18mm},
               level 3/.style={sibling distance=9mm},level 4/.style={sibling distance=6mm},
               edge from parent/.style={draw,thick,gray!60!black}},
  bk/.style={circle,draw=orange!70!black,fill=orange!22,minimum size=7.6mm,inner sep=0pt,font=\scriptsize,thick},
  bkm/.style={bk,fill=yellow!40,draw=orange!90!black,line width=1.2pt},
  ab/.style={rectangle,draw=gray!70,fill=gray!25,minimum size=2.2mm,inner sep=0pt},
  hoehe/.style={font=\scriptsize\bfseries,inner sep=1pt},
}
% Das Spielerverzeichnis (Mika, Emil, Tara, Bo, Ines, Pia, Zoe, Dana, Kai, Ruby, Vik); Argument: TikZ-Optionen
\newcommand{\verzeichnisbaum}[1][]{\begin{tikzpicture}[baum,#1]
\node[bk] (Mika) {Mika}
  child {node[bk] (Emil) {Emil}
    child {node[bk] (Bo) {Bo} child[missing] child {node[bk] (Dana) {Dana}}}
    child {node[bk] (Ines) {Ines} child[missing] child {node[bk] (Kai) {Kai}}}}
  child {node[bk] (Tara) {Tara}
    child {node[bk] (Pia) {Pia} child[missing] child {node[bk] (Ruby) {Ruby}}}
    child {node[bk] (Zoe) {Zoe} child {node[bk] (Vik) {Vik}} child[missing]}};
\end{tikzpicture}}

\begin{document}
\abtitel{2.4}{Höhe und Effizienz}

\textbf{Bot-Alarm!} 500 Bots haben sich alphabetisch registriert: Bot001, Bot002, …, Bot500. Seitdem ist der Login von Bot500 so langsam wie früher in der Liste. \web{Station 2.4, Baum-Architekt und Projekt Bot-Alarm}

\begin{aufgabe}[Baum-Architekt]{1}
\begin{lueckentext}
\textbf{a)} Mit den 7 Namen Ada, Ben, Cem, Dio, Eva, Fin, Gil ist die kleinste Höhe \lsg[2.25cm]{3}, zum Beispiel mit der Reihenfolge\lsgrest{Dio Ben Fin Ada Cem Eva Gil}
Die größte Höhe ist \lsg[2.25cm]{7}, zum Beispiel mit\lsgrest{Ada Ben Cem Dio Eva Fin Gil}
\textbf{b)} Für 15 Namen ist mindestens die Höhe \lsg[2.25cm]{4} nötig.
\end{lueckentext}\par\vspace{1mm}
\textbf{c)} So viele Knoten passen höchstens in einen Baum der Höhe h:\par\vspace{1mm}
{\renewcommand{\arraystretch}{1.1}
\begin{tabular}{|l|*{6}{>{\centering\arraybackslash}p{1.3cm}|}>{\centering\arraybackslash}p{2.2cm}|}\hline
Höhe h & 1 & 2 & 3 & 4 & 5 & 10 & h \\\hline
höchstens Knoten & \lsgtext{1} & \lsgtext{3} & \lsgtext{7} & \lsgtext{15} & \lsgtext{31} & \lsgtext{1023} & \schreibzeile\lsgtext{$2^h - 1$} \\\hline
\end{tabular}}
\end{aufgabe}

\begin{merke}[Merke: Höhe und Aufwand der Suche]
\begin{lueckentext}
Eine Suche besucht auf jeder Ebene höchstens \lsg[2.25cm]{einen} Knoten. Also gilt: Vergleiche im schlimmsten Fall $=$ \lsg[5.25cm]{Höhe}.\par
\textbf{Ausgewogen:} Alle Ebenen bis auf die letzte sind voll. Für n Knoten ist die Höhe etwa \lsg[3.5cm]{$\log_2(n+1)$}, aufgerundet. Die Suche wächst \lsg[5.25cm]{logarithmisch}: Doppelt so viele Accounts kosten nur \lsg[3.5cm]{einen} Vergleich mehr.\par
\textbf{Entartet:} Jeder Knoten hat höchstens ein Kind, der Baum ist eine \lsg[3.5cm]{Kette} mit Höhe n. Die Suche wächst \lsg[5.25cm]{linear} wie in der Liste. Das passiert zum Beispiel, wenn die Namen \lsg[3.5cm]{sortiert} eingefügt werden.\par
Die Form des Baums hängt nur von der \lsg[5.25cm]{Reihenfolge} beim Einfügen ab.
\end{lueckentext}
\end{merke}

\begin{aufgabe}[Wie viele Vergleiche im schlimmsten Fall? (Rechner auf der Website)]{2}
{\renewcommand{\arraystretch}{1.1}
\begin{tabularx}{\linewidth}{|l|*{4}{>{\centering\arraybackslash}X|}}\hline
Anzahl Accounts n & 15 & 1000 & 1 000 000 & 2 000 000 \\\hline
Liste & \lsgtext{15} & \lsgtext{1000} & \lsgtext{1 000 000} & \schreibzeile\lsgtext{2 000 000} \\\hline
entarteter Baum & \lsgtext{15} & \lsgtext{1000} & \lsgtext{1 000 000} & \schreibzeile\lsgtext{2 000 000} \\\hline
ausgewogener Baum & \lsgtext{4} & \lsgtext{10} & \lsgtext{20} & \schreibzeile\lsgtext{21} \\\hline
\end{tabularx}}
\end{aufgabe}

\newpage
\begin{aufgabe}[Die Höhe rekursiv berechnen (TODO 3)]{3}
\begin{minipage}[t]{0.52\linewidth}
\vspace{0pt}
\textbf{a)} Schreibe in jedes Kästchen die Höhe des Teilbaums, der an diesem Knoten beginnt. Fang unten an!\par\vspace{1mm}
\centering
\begin{tikzpicture}[baum,scale=1.05]
\node[bk] (Mika) {Mika}
  child {node[bk] (Emil) {Emil}
    child {node[bk] (Bo) {Bo} child[missing] child {node[bk] (Dana) {Dana}}}
    child {node[bk] (Ines) {Ines} child[missing] child {node[bk] (Kai) {Kai}}}}
  child {node[bk] (Tara) {Tara}
    child {node[bk] (Pia) {Pia} child[missing] child {node[bk] (Ruby) {Ruby}}}
    child {node[bk] (Zoe) {Zoe} child {node[bk] (Vik) {Vik}} child[missing]}};
\foreach \n/\h in {Mika/4,Emil/3,Tara/3,Bo/2,Ines/2,Pia/2,Zoe/2,Dana/1,Kai/1,Ruby/1,Vik/1}
  \node[draw,fill=white,minimum size=3.6mm,inner sep=0pt,font=\scriptsize\bfseries,anchor=south west] at ([xshift=-0.8mm,yshift=-0.8mm]\n.north east) {\lsgtext{\h}};
\end{tikzpicture}
\end{minipage}\hfill
\begin{minipage}[t]{0.45\linewidth}
\vspace{0pt}
\begin{lueckentext}
Dabei rechnet man von unten nach oben, also in \lsg[3.5cm]{Postorder}.\par
\end{lueckentext}
\textbf{b)} Ergänze den Code.\par
\begin{javaluecke}[basicstyle=\linespread{1.5}\footnotesize\ttfamily]
// Knoten
public int hoehe() {
   int l = §\lsg[3.2cm]{links.hoehe()}§;
   int r = §\lsg[3.2cm]{rechts.hoehe()}§;
   return §\lsg[3.8cm]{1 + Math.max(l, r)}§;
}
// Abschluss
public int hoehe() {
   return §\lsg[1.8cm]{0}§;
}
\end{javaluecke}
\end{minipage}
\end{aufgabe}

\begin{aufgabe}[Die Bot-Welle nachspielen (Projekt)]{4}
\begin{minipage}[t]{0.38\linewidth}
\vspace{0pt}
\textbf{a)} Skizze: Bot001 bis Bot004, alphabetisch eingefügt.\par\vspace{1mm}
\lsgkarobox{6cm}{2.5cm}{
  \foreach \n/\i in {001/0,002/1,003/2,004/3}{\node[bk,minimum size=5.5mm,font=\tiny,draw=red!75!black,fill=white] (b\i) at (0.6+1.45*\i,-0.4-0.55*\i) {\n};}
  \foreach \i [evaluate=\i as \j using int(\i+1)] in {0,1,2}{\draw[thick] (b\i) -- (b\j);}}
\end{minipage}\hfill
\begin{minipage}[t]{0.59\linewidth}
\vspace{0pt}\lueckentext
\textbf{b)} Bot-Welle: Höhe \lsg[2.25cm]{500}, Login Bot500: \lsg[2.25cm]{500} Vergleiche.\par
Grund:\lsgrest{Jeder neue Bot ist größer als alle bisherigen}
\lsgrest{und hängt ganz rechts: eine Kette.}
\end{minipage}
\end{aufgabe}

\begin{aufgabe}[Dieselben Bots in zufälliger Reihenfolge (Projekt)]{5}
\begin{lueckentext}
Drei Läufe: Höhe \lsg[1.5cm]{18}, \lsg[1.5cm]{20}, \lsg[1.5cm]{17}. Bestmöglich wäre für 500 Knoten die Höhe \lsg[2.25cm]{9}. Vergleich:\lsgrest{nicht perfekt, aber weit weg von 500.}
\end{lueckentext}
\end{aufgabe}

\begin{aufgabe}[Selbst bauen: ausgewogen einfügen (TODO 6)]{6}
\begin{minipage}[t]{0.42\linewidth}
\vspace{0pt}
\textbf{Vorgehensweise} für ein sortiertes Array (Bereich \code{von} bis \code{bis}):\par
\nr{1} Bereich leer? Dann fertig.\\
\nr{2} das mittlere Element einfügen\\
\nr{3} die linke Hälfte genauso\\
\nr{4} die rechte Hälfte genauso\par\vspace{1mm}
\begin{lueckentext}
\textbf{a)} Ada … Gil werden so eingefügt:\par\lsgrest{Dio Ben Ada Cem Fin Eva Gil}
Das ist die \lsg[3cm]{Preorder} des fertigen Baums.
\end{lueckentext}
\end{minipage}\hfill
\begin{minipage}[t]{0.55\linewidth}
\vspace{0pt}
\textbf{b)} Ergänze den Code.\par
\begin{javaluecke}[basicstyle=\linespread{1.5}\footnotesize\ttfamily]
public void ausgewogenEinfuegen(
        String[] sortiert, int von, int bis) {
   if (§\lsg[2.4cm]{von > bis}§) {           // §\nr{1}§
      return;
   }
   int mitte = §\lsg[3cm]{(von + bis) / 2}§;
   einfuegen(new Spieler(sortiert[mitte], 1));  // §\nr{2}§
   ausgewogenEinfuegen(sortiert, §\lsg[2.5cm]{von, mitte - 1}§); //§\nr{3}§
   ausgewogenEinfuegen(sortiert, §\lsg[2.5cm]{mitte + 1, bis}§); //§\nr{4}§
}
\end{javaluecke}
\end{minipage}
\end{aufgabe}
\end{document}
